\documentclass[12pt]{article}

% loads config.sty
\makeatletter
\def\input@path{{../../../tex/}}
\makeatother
\usepackage{config}
\usepackage{macros}

% =================================================================================================
% STYLE GUIDE -- for agents and collaborators editing this file. Not rendered; readers never see
% it. Keep it up to date when a convention changes.
% =================================================================================================
%
% SYMBOLS
%   y, g, e         phenotype, genetic (breeding) value, environmental component; y = g + e.
%   V_Y, V_A, V_E   their variances. Note V_Y, not V_P: the subscript names the VARIABLE whose
%                   variance it is, and the phenotype is always y.
%   h^2             heritability, V_A / V_Y.
%   x_{ik}, z_{ik}  standardized genotype of individual i at variant k, and its two per-allele
%                   halves; a (m) or (p) superscript says maternally or paternally inherited.
%   beta_k          the effect of variant k.
%   rho_y, rho_g    the correlation between mates' phenotypes and between their genetic values.
%                   rho_y is the model's input and is fixed across generations; rho_g is induced,
%                   and at equilibrium rho_g = h^2 rho_y.
%   mu              the RAW co-path coefficient, rho_y / V_Y. A co-path is labelled in figures
%                   with the correlation it induces, [rho_y], and mu is left to be derived.
%
% SUPERSCRIPTS AND SUBSCRIPTS
%   A PARENTHESIZED superscript is a generation or a step count: V_A^{(0)} is the base
%     population, V_A^{(t)} generation t, V_A^{(eq)} equilibrium; \Mu^{(n)} was n mating steps.
%     Macros exist for the common ones -- \VAo, \VYo, \VAt, \VYt, \VAeq, \VYeq, \hsqo, \hsqt.
%   A BARE superscript is an iteration count, not a generation.
%   From Section 2.3 onward everything is at equilibrium and the (eq) superscript is DROPPED;
%     that is stated once at the top of 2.3.1 and then relied on.
%   A \perp superscript marks a RESIDUAL -- the part of a quantity orthogonal to its expectation
%     -- and the base symbol says what kind of thing it is. So z^{\perp(m)}_{ok} is the part of
%     the maternally inherited allele not predicted by the mother's genotype, and g^{\perp}_o the
%     part of the offspring's genetic value not predicted by the parental mean. Do NOT use
%     \epsilon or \varepsilon for these; \epsilon is reserved for path-analysis DISTURBANCES.
%
% GRAPH OBJECTS (Section 2.3)
%   One node set -- individuals -- carrying two edge sets. Undirected edges are matings and give
%   the MATING GRAPH; directed parent-to-offspring edges are meioses and give the MEIOSIS GRAPH.
%   Calligraphic letters name objects in them, with a matching N for the distance or count:
%     \Mu  = \mathcal{M}   mating chain,   \Nmu  = N_\mathcal{M}   mating distance
%     \Lin = \mathcal{L}   lineal chain,   \Nlin = N_\mathcal{L}   lineal distance
%     \Anc = \mathcal{A}   ancestors
%     \mathcal{P}          a path,         N_o                     out-degree in the chain digraph
%   Distances are infinity, not undefined, when no connection exists -- that makes the
%   dominant-path comparison work without a case split.
%
% MACROS
%   All live in ../../../tex/macros.sty; boxes and markup in ../../../tex/config.sty. Prefer
%   adding a macro to repeating a construction. Two traps:
%     - NEVER write \mathbb{1}. msbm has no digit glyphs, so it silently renders an unrelated
%       symbol with no warning. Use \mathbf{1} (there is an \IMu macro for the mating indicator).
%     - Use \overline and \widehat, not \bar and \hat, over calligraphic letters; the short
%       accents sit off-centre.
%   \add / \remove and the added / removed environments are REVIEW SCAFFOLDING for showing a
%   rewrite against what it replaces. Strip them once a revision is accepted.
%
% WRITING PREFERENCES
%   - Cut fluff. Expand a derivation where the expansion earns its space, not by default.
%   - Lean on a small number of variables. Do not introduce a symbol that is used briefly and
%     then dropped; if a quantity appears once, write it out instead.
%   - Define a technical term the first time it is used, in one clause, in plain words
%     ("a graph with no cycles is a forest"). Do not prove standard properties of standard
%     objects -- name the object and rely on it.
%   - Avoid very compact terminology unless it is well explained.
%   - "lineal", not "linear", for descent relationships.
%   - Assumption boxes answer "what must I believe?"; approximation boxes answer "what am I
%     neglecting?" and carry a \justification. Every assumption gets a STABLE id, reused verbatim
%     across writeups, and is cited with \assumptionid / \approximationid.
%
% FIGURES AND CHECKS
%   Every .tikz in figures/ is GENERATED by figures/make_figures.py, which also asserts every
%   hand-typed table entry and boxed equation against pathMgr. Run
%       python figures/make_figures.py            # regenerate the figures and run the checks
%       python figures/make_figures.py --check    # checks only, write nothing
%   before committing. A hand-typed result that is not asserted there is a result that will drift.
%   Regenerating must leave the .tikz files byte-identical unless a figure actually changed.
% =================================================================================================

\title{Relative Covariance}
\date{\today}

\begin{document}

\maketitle

\textit{\textbf{Note:}
Parts of this document may be AI-written.
My general workflow is to write sections of the article myself, laying out the general points I want to hit on and going into detail where relevant.
I then converse with an LLM to identify errors in the logic or math and opportunities to use clearer or more generalizable notation.
I then have the LLM make the edits to reflect our proposed changes, including adding more detailed derivations.}

This article describes how the phenotypic covariance between two individuals $i$ and $j$ can be modeled by limited information we have about the pair.
I derive formulas from the starting trait generative model and only introduce assumptions at the point the derivation requires it.
I start with a simple trait architecture and iteratively add more complexity.

\section{Additive genetics under random mating with independent environment}\label{sec:random-mating}
We begin with the simplest architecture: a trait determined by an additive genetic component and an independent environmental component, in a randomly mating population.

\subsection{The phenotype model}
For an individual $i$ we write the (standardized) phenotype as
\begin{equation}\label{eq:model}
    y_i = g_i + e_i ,
\end{equation}
where $g_i$ is the genetic value and $e_i$ the environmental deviation.
We take $y$ to be standardized across the population, $\E[y]=0$ and $\Var[y]=1$, and both components to be mean-centered, $\E[g]=\E[e]=0$, with variances $\Var[g]=V_A$ and $\Var[e]=V_E$.

We then introduce the following assumptions.

\begin{assumption}{Additivity}{Additive genetic value}
The genetic value is a fixed additive function of the individual's own genotype,
$$ g_i = \sum_{k=1}^{M} \beta_k\, x_{ik} , $$
where $\beta_k$ is the per-standardized-allele effect of variant $k$, $x_{ik}$ is the standardized genotype ($\E[x_{ik}]=0$, $\Var[x_{ik}]=1$), and $M$ is the number of variants in the genome.
Non-causal variants are included with $\beta_k=0$.
Because we assume an additive model, the genetic value can also be called a \emph{breeding value}, although we stick to the former term.
\bundles{the additive model; no dominance, no gene--gene interactions (GxG, epistasis), and no gene--environment interactions (GxE) --- each would make the genetic value depend on more than the individual's own genotype in an additive way.}
\end{assumption}

\begin{assumption}{GE-indep}{Gene--environment independence}
An individual's environmental component is independent of their own genetic component:
$$ g_i \perp e_i $$
\bundles{no gene--environment correlation or interaction, within an individual.}
\end{assumption}

\begin{assumption}{No-IGE}{No indirect genetic effects}
One individual's genetic component is independent of another individual's environmental component:
$$ g_i \perp e_j \ \ (i \neq j) . $$
\bundles{no indirect (e.g.\ parental) genetic effects --- one individual's genotype does not shape another's environment.}
\end{assumption}

\begin{assumption}{Env-indep}{Independent environments}
The environmental component is independent across individuals, including relatives:
$$  e_i \perp e_j \ \ (i \neq j). $$
\bundles{no shared/common environment ($C$); no environmental confounder shared by the pair (e.g.\ population stratification).}
\end{assumption}

Because $g_i\perp e_i$ (\assumptionid{GE-indep}), $\Cov[g_i,e_i]=0$ and the phenotypic variance decomposes as $\Var[y]=V_A+V_E$.
Narrow--sense heritability is typically defined as:
$$ h^2 = \frac{V_A}{\Var[y]}$$
Since $\Var[y] = 1$, $h^2 = V_A$ under this model.


\subsection{Phenotypic covariance is genotypic covariance}
For a single fixed pair, $y_i$ and $y_j$ are two fixed values, so ``their covariance'' only has meaning as an average over an ensemble of pairs.
The ensemble we use is the set of all pairs that share a given piece of \emph{relatedness information} $\mathcal{I}_{ij}$ about the pair, and the object of interest is the conditional cross-product
\begin{equation}\label{eq:object}
    C_{ij}(\mathcal{I}) \;:=\; \E[\, y_i y_j \mid \mathcal{I}_{ij}\,] .
\end{equation}

If $\E[y_i \mid \mathcal{I}]=0$, then $C_{ij}(\mathcal{I}) = \Cov[y_i,y_j\mid \mathcal{I}]$.
For most information $\mathcal{I}$, like the realized relatedness or degree between two relatives, this equivalence is applicable.
An exception is when the information we are modelling contains individuals' actual genotype information, for which $\E[y_i\mid \text{genotypes}]=g_i\neq 0$, therefore the expected cross-product of two individuals is not equivalent to their covariance given the fixed information.

Using \eqref{eq:model} together with \assumptionid{No-IGE} and \assumptionid{Env-indep}, every cross term involving $e$ vanishes in expectation (each factors through a mean-zero environmental term), so the phenotypic cross-product reduces to the genetic one:
\begin{equation}\label{eq:reduce}
    C_{ij}(\mathcal{I}) \;=\; \E[\,(g_i+e_i)(g_j+e_j) \mid \mathcal{I}\,] \;=\; \E[\,g_i g_j \mid \mathcal{I}\,] .
\end{equation}
The rest of the article evaluates \eqref{eq:reduce} at four levels of resolution for $\mathcal{I}$.

\subsection{Level 0 --- the exact genetic cross-product}
At the finest resolution we condition on the realized genotypes of $i$ and $j$ at the causal variants, for which we know the true causal effect sizes.
Then $g_i$ and $g_j$ are fixed, and by \assumptionid{Additivity},
\begin{equation}\label{eq:level0}
    \boxed{ C_{ij} \;=\; \E[\, y_i y_j \mid \bx_i, \bx_j, \bbeta \,] \;=\; g_i g_j \;=\; \sum_{k}\sum_{k'} \beta_k \beta_{k'}\, x_{ik}\, x_{jk'} }
\end{equation}
Under our model and assumptions so far, this returns the true expected cross-product of two individuals' phenotypes.
The genetic cross-product between two individuals is, fundamentally, an effect-size-weighted function of the genotypes they actually carry at causal sites.
Notably, this formula does not require any knowledge of the relatedness of the two individuals.

\subsection{Level 1 --- genome-wide allelic sharing}
If we do not know the causal effect sizes, then we cannot apply \eqref{eq:level0}.
However, we can assume that genetic effects' cross-products are independent of individuals' allelic sharing patterns.

\begin{assumption}{Random-effects}{Random effects, exchangeable across variants}
Treating the effects as random, they are mean-zero and uncorrelated across variants, the distribution of their cross-products $\{\beta_k \beta_{k'}\}$ is independent of the genotypes $\bX$, and every variant carries the same second moment:
$$ \E[\beta_k]=0, \quad \E[\beta_k \beta_{k'}] = 0 \ \ (k \neq k'), \quad \E[\beta_k \beta_{k'} \mid \bX] = \E[\beta_k \beta_{k'}], \quad \E[\beta_k^2] = \E[\beta_{k'}^2] . $$
The last condition is a statement about the \emph{prior} we place on effects, not about their realized values: it says we have no information distinguishing one variant from another, so we weight them symmetrically.
\bundles{the infinitesimal / random-effects model.}
\end{assumption}

Taking the expectation of \eqref{eq:level0} over the effects, and using the independence of the cross-products from the genotypes (\assumptionid{Random-effects}) to replace $\E[\beta_k \beta_{k'} \mid \bX]$ by $\E[\beta_k \beta_{k'}]$:
$$
    \E[\,g_i g_j \mid \bx_i, \bx_j \,]
    = \sum_{k}\sum_{k'} \E[\beta_k \beta_{k'}]\, x_{ik} x_{jk'}.
$$
Because the effects are mean-zero and uncorrelated across variants, $\E[\beta_k\beta_{k'}] = 0$ whenever $k\neq k'$; the entire off-diagonal ($k\neq k'$) collapses regardless of any correlation between the genotypes.
What remains is the diagonal, weighted by each variant's effect variance:
$$
    \E[\,g_i g_j \mid \bx_i, \bx_j \,]
    = \sum_{k} \E[\beta_k^2]\, x_{ik} x_{jk}.
$$
The additive variance is the total of these: since $\Var[x_{ik}]=1$ and the off-diagonal effect terms vanish, $V_A = \Var[g_i] = \sum_k \E[\beta_k^2]$, and thus $\E[\beta_k^2] = V_A/M$.
Effectively, the per-variant effect weighting of \eqref{eq:level0} has been replaced by a uniform weighting across all variants.
Therefore
\begin{equation}
    C_{ij}(\bx_i, \bx_j) \;=\; V_A \cdot \frac{1}{M}\sum_{k} x_{ik} x_{jk} .
\end{equation}
The quantity multiplying $V_A$ is the genome--wide genotype cross-product (i.e. allelic sharing), which we'll define by:
$$A_{ij} =  \frac{1}{M}\sum_{k} x_{ik} x_{jk}$$
This can be thought of as an identity-by-state (IBS)-based measure of ``relatedness'', and is exactly the $(i,j)$ entry of what is typically computed as a Genetic Relatedness Matrix (GRM), $\bA$.
Therefore, we can further simplify:
\begin{equation}\label{eq:level1}
    \boxed{ C_{ij}(A_{ij}) \;=\; V_A A_{ij} }
\end{equation}

\subsection{Level 2 --- realized relatedness}
The problem with calculating $A_{ij}$ is that we may not have genotyped every causal variant that contributes to a trait.
As long as some of the causal variants are not in perfect correlation with a genotyped variant, then the covariance of two individuals' phenotypes will not fully take into account the contributions of the ungenotyped causal variants.
However, when individuals are related, their IBS-based relatedness from genotyped variants is informative of the IBS-based relatedness of ungenotyped variants.
Therefore, the former can be used to capture the contributions of the latter.

When individuals inherit a genomic segment from a common ancestor, then all of the variants within that genomic segment will be identical-by-state, as long as we assume that de-novo mutations have not caused the inherited alleles to be different.
These genomic segments and their variants are said to be identical-by-descent (IBD).

\begin{assumption}{No-mutation}{IBD implies IBS}
Alleles that are identical-by-descent are also identical-by-state: descent is transmitted without mutation.
\bundles{the infinite-sites / no-recurrent-mutation assumption.}
\end{assumption}

When two individuals' alleles are IBD, then they are always IBS; when they are not IBD, they can still be IBS.
Therefore, we can decompose the IBS-based relatedness $A_{ij}$ into a sum over IBD allele pairs and a sum over non-IBD allele pairs.
Write each standardized genotype as the sum of its two per-allele (maternal and paternal) standardized values, $x_{ik} = z^{(m)}_{ik} + z^{(p)}_{ik}$.
Since $\Var[x_{ik}]=1$ and the two alleles are independent (\assumptionid{No-inbreeding}), each per-allele value has $\E[z]=0$ and $\Var[z]=\tfrac12$.
The genotype cross-product then expands over the four ordered allele pairs, and $A_{ij}$ splits by IBD status:
\begin{equation}\label{eq:Adecomp}
\begin{split}
    A_{ij}
    &= \frac{1}{M}\sum_{k} x_{ik} x_{jk}
    = \frac{1}{M}\sum_{k} \sum_{u,v\in\{m,p\}} z^{(u)}_{ik} z^{(v)}_{jk} \\
    &= \underbrace{\frac{1}{M}\sum_{\text{IBD pairs}} z^{(u)}_{ik} z^{(v)}_{jk}}_{\text{IBD}}
    \;+\; \underbrace{\frac{1}{M}\sum_{\text{non-IBD pairs}} z^{(u)}_{ik} z^{(v)}_{jk}}_{\text{non-IBD}} .
\end{split}
\end{equation}

For an IBD allele pair, the two per-allele values are copies of the same ancestral allele and, by \assumptionid{No-mutation}, of the same state; they are therefore equal, so their expected cross-product is $\E[(z^{(u)}_{ik})^2] = \Var[z^{(u)}_{ik}] = \tfrac12$.
For non-IBD allele pairs, we have to make an assumption on how they are distributed.

\begin{assumption}{Non-IBD-indep}{Independence of non-IBD alleles}
Any two alleles that are \emph{not} IBD are independent draws from the population allele distribution.
\bundles{between individuals: random mating, no population structure, and no assortative mating.}
\end{assumption}

\begin{assumption}{No-inbreeding}{No inbreeding}
The maternally and paternally inherited alleles of an individual are not IBD with each other.
\bundles{a randomly mating, outbred population (Hardy--Weinberg within individuals).}
\end{assumption}

Two consequences follow, and both are used throughout.
First, an individual's own two alleles, not being IBD, are independent draws from the population by \assumptionid{Non-IBD-indep}, and so are uncorrelated with each other.
A variant's genotype variance is therefore the sum of its two allele variances, $2p_k(1-p_k)$, which is what makes the standardized genotype $x_{ik}$ have variance $1$ and each of its two per-allele values variance $\tfrac12$.
Second, IBD is an equivalence relation on alleles --- two are IBD exactly when they are copies of the same ancestral allele --- and the assumption says no class may contain both alleles of one individual.
At most two of the four ordered allele pairs between a pair of individuals can therefore be IBD, which is what keeps the number of alleles they share IBD at a variant in $\{0,1,2\}$.

By the definition of independent draws, the expected cross-product of a non-IBD allele pair is therefore $\E[z^{(u)}_{ik}]\,\E[z^{(v)}_{jk}] = 0$.
Letting $\IBD_{ij}(k)\in\{0,1,2\}$ denote the number of alleles the pair shares IBD at variant $k$, exactly $\IBD_{ij}(k)$ of the four ordered allele pairs are IBD --- each contributing $\tfrac12$ --- and the rest contribute $0$, so
$$
    \E\!\left[\,x_{ik} x_{jk} \,\middle|\, \IBD_{ij}(k)\,\right]
    = \tfrac12 \IBD_{ij}(k) .
$$
A variant thus contributes expected cross-product $0$, $\tfrac12$, or $1$ according to whether the pair shares $0$, $1$, or $2$ alleles IBD there.
We define the pair's \emph{realized relatedness} $\pi_{ij}$ as the average IBD state across variants, halved --- equivalently, the fraction of the pair's genome that is IBD:
\begin{equation}\label{eq:pi-def}
    \pi_{ij} \;:=\; \frac{1}{2M}\sum_{k} \IBD_{ij}(k) .
\end{equation}

Conditioning on the pair's realized IBD states and using the per-variant identity $\E[x_{ik}x_{jk}\mid \IBD_{ij}(k)] = \tfrac12\,\IBD_{ij}(k)$, linearity of expectation gives the expected allelic sharing:
$$
    \E\!\left[\,A_{ij}\,\middle|\,\{\IBD_{ij}(k)\}\,\right]
    = \frac{1}{M}\sum_{k} \tfrac12\,\IBD_{ij}(k)
    = \frac{1}{2M}\sum_{k} \IBD_{ij}(k)
    = \pi_{ij} .
$$
This depends on the IBD states only through $\pi_{ij}$, so we may condition on $\pi_{ij}$ alone.
Crucially, it required no assumption that the per-variant IBD states are independent --- and indeed they are not: an IBD state changes only at a cross-over within a co-inherited segment, so physically linked variants have strongly correlated IBD states.
Linearity of expectation is insensitive to this correlation.

Combining with \eqref{eq:level1}, the pair's expected phenotypic cross-product is their realized relatedness scaled by the additive variance:
\begin{equation}\label{eq:level2}
    \boxed{\ \E[\,y_i y_j \mid \pi_{ij}\,] \;=\; V_A\,\E[A_{ij}\mid\pi_{ij}] \;=\; V_A\,\pi_{ij}\ } .
\end{equation}

\subsection{Level 3 --- genetic degree}
In cases where we cannot measure $\pi_{ij}$ based on per-variant IBD states, we can still estimate it if we know the genetic degree $d_{ij}$ between the individuals.

\begin{assumption}{Genetic-degree}{Degree determines expected relatedness}
The genetic degree $d_{ij}$ fixes the distribution of IBD sharing through Mendelian segregation, so that the expected relatedness given the degree is a single number $\bar\pi_{d_{ij}}$:
$$ \E[\pi_{ij}\mid d_{ij}] = \bar\pi_{d_{ij}} . $$
\end{assumption}

Furthermore, under the outbreeding assumption (\assumptionid{No-inbreeding}) --- applied here not only to $i$ and $j$ but to the entire pedigree connecting them, so that it contains no inbreeding loops --- $\bar\pi_{d_{ij}} = 2^{-d_{ij}}$, which is not fully derived here.
We can then obtain the expected phenotypic covariance given the degree $d_{ij}$:
\begin{equation}\label{eq:level3}
    \E[\,y_i y_j \mid d_{ij}\,] \;=\; \E\big[\,\E[y_i y_j\mid \pi_{ij}]\,\big|\, d_{ij}\,\big] \;=\; V_A\, \E[\pi_{ij}\mid d_{ij}] \;=\; V_A\,\bar\pi_{d_{ij}} .
\end{equation}

\subsection{Summary}
Each level conditions on coarser information than the one above and is its conditional expectation; the rightmost column lists the assumption(s) newly required to descend one step.

Explicitly, each level is the conditional expectation of the finer level above it, averaging out one additional layer of information:
\begin{align*}
    \E[y_i y_j \mid \bx_i,\bx_j]
      &= \E\big[\,g_i g_j \,\big|\, \bx_i,\bx_j\,\big] = V_A\,A_{ij}, \\
    \E[y_i y_j \mid \pi_{ij}]
      &= \E\big[\,V_A\,A_{ij} \,\big|\, \pi_{ij}\,\big] = V_A\,\pi_{ij}, \\
    \E[y_i y_j \mid d_{ij}]
      &= \E\big[\,V_A\,\pi_{ij} \,\big|\, d_{ij}\,\big] = V_A\,\bar\pi_{d_{ij}},
\end{align*}
starting from the Level-0 value $\E[y_iy_j\mid\bx_i,\bx_j,\bbeta]=g_ig_j$.
Equivalently, the ladder telescopes into a single nested expectation,
$$
    \E[y_i y_j \mid d_{ij}]
    = \E\Big[\,\E\big[\,\E[\,g_i g_j \mid \bx_i,\bx_j\,]\,\big|\,\pi_{ij}\,\big]\,\Big|\,d_{ij}\,\Big]
    = V_A\,\bar\pi_{d_{ij}} .
$$
This is valid down the generative chain $d_{ij}\to\pi_{ij}\to(\bx_i,\bx_j)\to y$: conditioned on the finer variable, the phenotype does not depend further on the coarser one, so the tower property applies at each step.

\begin{center}
\small
\begin{tabular}{c l l p{4.3cm}}
\textbf{Level} & \textbf{Conditioned on} & $\E[y_i y_j \mid \cdot]$ & \textbf{New assumptions} \\
\hline
0 & causal genotypes & $g_i g_j$ & \assumptionid{Additivity} (given \assumptionid{No-IGE}, \assumptionid{Env-indep}) \\
1 & realized sharing $A_{ij}$ & $V_A\,A_{ij}$ & \assumptionid{Random-effects} \\
2 & realized relatedness $\pi_{ij}$ & $V_A\,\pi_{ij}$ & \assumptionid{No-mutation}, \assumptionid{Non-IBD-indep}, \assumptionid{No-inbreeding} \\
3 & genetic degree $d_{ij}$ & $V_A\,\bar\pi_{d_{ij}}$ & \assumptionid{Genetic-degree} \\
\hline
\end{tabular}
\end{center}

The property that makes this a ladder is worth naming, because the next architecture breaks it.
At each step, the finer variable \emph{screens off} the coarser one: once we condition on $\pi_{ij}$, knowing $d_{ij}$ as well would tell us nothing more about $y_iy_j$, and likewise for the other steps.
Screening is what licenses the tower property, and it is the reason the four information sets can be arranged in a single descending chain rather than a more complicated structure.

\section{Additive genetics under assortative mating}\label{sec:am}
We will now do away with the assumption of random mating (particularly implied in \assumptionid{Non-IBD-indep}) and instead provide a model for assortative mating that will let us quantify its impact on phenotypic covariance.

The phenotype model remains the same: we assume an additive genetic architecture (\assumptionid{Additivity}) with an independent environmental component (\assumptionid{GE-indep}).
As a consequence of the assortative mating model we adopt below, covariance between the environmental and genetic components between some individuals will arise, breaking the weakest forms of the assumptions \assumptionid{No-IGE} and \assumptionid{Env-indep}.
We also still assume no new mutations (\assumptionid{No-mutation}) and no inbreeding (\assumptionid{No-inbreeding}).

\subsection{The mating model}
We assume a model of assortative mating based on mates' similarity in a sorting factor, which is also our focal trait.

\begin{assumption}{Primary-AM}{Primary Assortative Mating}
Mating is non-random such that two mates' phenotypes are correlated for some trait at the time of mating:
$$ \Corr[\,y_m, y_p\,] = \rho_y , $$
where the subscripts $m$ and $p$ index the maternal and paternal member of a mated pair.
Matching is on the phenotype \emph{only}: any covariance between the mates' trait components is a downstream consequence of this phenotypic matching, and is not specified separately.
Following the convention that a subscript names the component being correlated, we write $\rho_g$ for the induced correlation between mates' genetic values:
$ \Corr[\,g_m, g_p\,] = \rhogt $, where the superscript records that this induced correlation changes from one generation to the next, unlike $\rho_y$, which is held fixed across generations.
\bundles{primary phenotypic assortment, as opposed to social homogamy (matching on a shared environment) or genetic homogamy (matching on genotype directly); it also excludes convergence, where mates grow more similar after mating. $\rho_y$ is often written $\mu$ or $r$ in the literature.}
\end{assumption}

\subsection{Dynamics and equilibrium}
To make modelling the dynamics of populations easier, we will make the following assumption: 

\begin{assumption}{Discrete-generations}{Discrete, non-overlapping generations}
The population is partitioned into generations indexed by $t$, and every individual mates within their own generation.
\bundles{no overlapping generations and no age structure; in particular no mating between individuals of different generations, and no individual contributing offspring to more than one generation.}
\end{assumption}

When starting from a randomly mating population, one generation of assortative mating causes the immediate offspring generation to exhibit different phenotypic and genetic properties on expectation.
Luckily, if the parameters of assortative mating remain fixed, an equilibrium is reached such that stable properties of the population can be studied.
We will first explore how assortative mating affects statistical properties of the population during the first few generations, then we will solve for the parameter values that lead to an equilibrium state.

\subsubsection{Path Analysis}\label{sec:path-analysis}
First, we'll introduce path analysis, a very helpful tool for modelling the covariance between variables in a causal system, such as genetic inheritance and assortative mating.
Path analysis was introduced by Sewall Wright \citep{Wright1921JAgricRes}, as a way of decomposing an observed correlation between variables into the contributions of a hypothesized causal structure, and set out systematically thirteen years later as the method of path coefficients \citep{Wright1934AnnMathStat}.
Wright applied it to mating systems in the same period \citep{Wright1921GeneticsMatingI, Wright1921GeneticsMatingIII}; together with Fisher's treatment of the same problem \citep{Fisher1918TransRoySocEdin}, that work is the origin of the classical results on how assortment changes the correlation between relatives --- the results this section rederives.
The explanations below are adapted from the Supplementary Note of \citet{Sunde2025NatComm}, which introduced the co-path edge type (described below).
Unlike \citet{Sunde2025NatComm}, we do not assume that the variables are standardized, so we will keep track of variances explicitly.

A \emph{path diagram} is a graphical representation of a generative model.
Each variable is a node --- drawn as a rectangle if it is observed and an ellipse if it is latent --- and the nodes are joined by three kinds of edges:

\begin{itemize}
    \item A \textbf{directed} edge $a \rightarrow b$ labelled with a coefficient $c$ says that $a$ is a cause of $b$, contributing $c\,a$ to $b$'s structural equation.
    Collecting the arrows that point \emph{into} a node recovers that node's equation, so the diagram and the model equations carry the same content.
    \item A \textbf{bidirected} edge $a \leftrightarrow b$ carries covariance between the \emph{disturbances} of $a$ and $b$.
    A variable's disturbance is what remains of it once its causes in the diagram are accounted for: if the arrows into $b$ give $b = c_1 a_1 + \cdots + c_k a_k + \epsilon_b$, then $\epsilon_b$ is $b$'s disturbance, and for a variable with no incoming arrows the disturbance is the variable itself.
    The self-loop $a \leftrightarrow a$ is the case $b=a$: it is $\Var[\epsilon_a]$, the part of $\Var[a]$ not supplied by $a$'s causes, so a variable with no incoming arrows carries its entire variance on such a loop.
    \item A \textbf{co-path} $a \;\text{---}\; b$, drawn as a plain line with no arrowheads, is covariance attributable to \emph{matching}: $a$ and $b$ came to be associated because they were paired on their values, not because they share a cause.
\end{itemize}

The covariance between \emph{any} two variables in the diagram, latent or observed, can be calculated quickly by following Wright's tracing rules.
Enumerate every \emph{valid chain} joining the two variables; their covariance is the sum over chains of the product of the coefficients each chain passes through.
A valid chain between $a$ and $b$ has the shape
$$
    a \leftarrow \cdots \leftarrow u \;\leftrightarrow\; v \rightarrow \cdots \rightarrow b :
$$
trace \emph{backward} along directed edges from $a$ to some ancestor $u$, cross \emph{exactly one} bidirected edge $u \leftrightarrow v$ (possibly a self-loop, $u=v$), then trace \emph{forward} along directed edges from $v$ down to $b$.
A chain turns around exactly once, and it turns around on a bidirected edge, because that is the only place variance enters the system.

Co-paths extend this.
A chain may pass \emph{through} a co-path and keep going: a co-path joins two valid chains of the kind above into a longer one,
$$
    \underbrace{a \leftarrow \cdots \leftrightarrow \cdots \rightarrow y_m}_{\text{valid chain}}
    \;\text{---}\;
    \underbrace{y_p \leftarrow \cdots \leftrightarrow \cdots \rightarrow b}_{\text{valid chain}} .
$$
Two further rules apply \citep{Sunde2025NatComm}: a chain can neither start nor end with a co-path, and a chain cannot use two co-paths belonging to the same mating process.
It \emph{may} cross co-paths of different mating processes.

A bidirected edge is a statement about disturbances, so a chain terminates on it: a disturbance has no causes in the diagram, so there is nothing to keep tracing back through, and the chain cannot reach past $y_m$ to the things that caused $y_m$.
A co-path is a statement about the matched variables themselves, so a chain continues backward through it, which encodes the fact that matching two people on their phenotypes induces correlation between \emph{all the causes} of those phenotypes --- their genetic values, and every individual variant behind them.
This is the mathematical basis behind the \assumptionid{Primary-AM} assumption.

The quantity multiplied along a chain when it crosses a co-path is not the correlation $\rho_y$ itself, but
\begin{equation}\label{eq:copath-mu}
    \mut \;=\; \frac{\rho_y}{\sqrt{\Var[y^{(t)}_m]\,\Var[y^{(t)}_p]}} \;=\; \frac{\rho_y}{\VYt} ,
\end{equation}
because a co-path with coefficient $\mu$ contributes $\Cov[a,b] = \mu\,\Var[a]\,\Var[b]$ \citep{Sunde2025NatComm}.
The two coincide only when the matched variables have unit variance.
Note the use of a superscript to denote the generation a quantity belongs to, since $\VYt$ changes from one generation to the next under assortative mating (described below).
A value of $\mu$ carried unchanged from one generation to the next is therefore wrong, which is the practical reason we label a co-path with its correlation in brackets, as $[\rho_y]$, and leave $\mut$ to be derived from the diagram's own variances.

Tracing chains by hand is instructive for a single pair and hopeless for a pedigree, since the number of chains grows multiplicatively with depth.
The diagrams and covariances in this section were produced with \texttt{pathMgr}, a small package written alongside this writeup: it takes a model specification, returns the covariance between any two of its variables as an algebraic expression, and emits the path diagram to go with it.
It computes each covariance in two independent ways --- by enumerating chains exactly as above, and in closed form from the matrix identity $\bSigma = \B{F}(\bI-\bA)^{-1}\bS(\bI-\bA)^{-\intercal}\B{F}^{\intercal}$ --- and their agreement is the correctness check.

That identity is the diagram written as matrices, one row and column per variable.
Collecting all the variables into a vector $\bv$, the matrix $\bA$ holds the \textbf{directed} edges --- $A_{ij}$ is the coefficient on the arrow from $j$ to $i$ --- so the whole system of structural equations is $\bv = \bA\bv + \bepsilon$, where $\bepsilon$ is the vector of disturbances.
The symmetric matrix $\bS$ holds the \textbf{bidirected} edges: $S_{ij}$ is the value of $i \leftrightarrow j$, and its diagonal $S_{ii}$ is the self-loop, the variance of $i$'s disturbance.
$\B{F}$ merely selects rows, keeping the variables we want reported and discarding the rest.

The inverse is what does the tracing.
In an acyclic model $\bA$ is nilpotent, so $(\bI-\bA)^{-1} = \bI + \bA + \bA^2 + \cdots$, and the $(i,u)$ entry of that sum is the total over all directed paths from $u$ to $i$ of the product of their coefficients --- one term per path, $\bA^n$ contributing the paths of length $n$.
Sandwiching $\bS$ between it and its transpose therefore gives, for latent and observed variables alike,
$$
    \Sigma_{ij} \;=\; \sum_{u}\sum_{v}
    \underbrace{\big[(\bI-\bA)^{-1}\big]_{iu}}_{i \leftarrow \cdots \leftarrow u}
    \underbrace{S_{uv}}_{u \leftrightarrow v}
    \underbrace{\big[(\bI-\bA)^{-1}\big]_{jv}}_{v \rightarrow \cdots \rightarrow j} ,
$$
which is Wright's rule term for term: a backward leg, exactly one bidirected edge, a forward leg, summed over every chain, indexed by the $u \leftrightarrow v$ it turns around on.
The two methods are the same statement --- one enumerated, one summed in closed form --- which is why their agreement tests the implementation rather than the mathematics.
Co-paths are the one thing that sits outside this identity and contributes further terms; there, the chain rules given above are the definition.

\subsubsection{Consequences of assortative mating on first generation of mates}
We assume a population has been randomly mating as described by assumption \assumptionid{Non-IBD-indep} up to some generation $t=0$, which can also be referred to as the base or randomly-mating population.
The parameters defined in the previous section, such as the additive variance $V_A$ and heritability $h^2$, are measured in this base population and denoted by a $(0)$, e.g. $\VAo$ and $\hsqo$, since these terms do not remain fixed across generations under assortative mating.
Note that $V_E$ remains fixed.

\subsubsection*{Building the model from haploid alleles}
To properly model the effects of assortative mating, it pays off to think in terms of individual haploid alleles, rather than diploid genotypes.
We previously defined $z_{ik}^{(m)}$ and $z_{ik}^{(p)}$ as the standardized values of the maternally- and paternally-inherited alleles at variant $k$ for individual $i$, and $x_{ik} = z_{ik}^{(m)} + z_{ik}^{(p)}$ as the standardized diploid genotype.
Each haploid allele $a^{(m)}$ and $a^{(p)}$ can take on two values, $0$ or $1$, and therefore each has a variance of $p_k(1-p_k)$, where $p_k$ is the allele frequency at variant $k$.
Typically, diploid genotypes are standardized by their expected variance when the two alleles are independent, which is $2p_k(1-p_k)$, resulting in a standardized variance of 1, as used in the previous section.
Under this standardization, each haploid allele has a standardized variance of $\tfrac12$, and if we assume that the allele frequency does not change over time, then neither does this quanity:

\begin{assumption}{Constant-freq}{Allele frequencies do not change}
The allele frequency $p_k$ of each variant is the same across every generation.
As a consequence, the standardization divides by a fixed quantity, and the standardized variance of each haploid allele is the same in every generation:
$$
    z_{ik} \;=\; \frac{a_{ik} - p_k}{\sqrt{2p_k(1-p_k)}}
    \qquad\Longrightarrow\qquad
    \Var[z_{ik}] \;=\; \frac{\Var[a_{ik}]}{2p_k(1-p_k)} \;=\; \frac{p_k(1-p_k)}{2p_k(1-p_k)} \;=\; \tfrac12 ,
$$
where $a_{ik} \in \{0,1\}$ is the unstandardized haploid allele count and $p_k$ is the allele frequency at variant $k$.
\bundles{no selection on the trait, no drift, no migration; in a finite population it is the large-$N_e$ limit \citep{Yengo2018NatHumBehav}.}
\end{assumption}

In contrast, the standardized diploid genotype can have non-unit variance and change over time if the two alleles are correlated, as will be demonstrated in the case of assortative mating.
Our path diagram will therefore start with haploid alleles which are summed together to form diploid genotypes, which are then summed to form the additive genetic value, which is then added to the environmental value to form the phenotype.

\subsubsection*{The diagram}
Mates in generation $t=0$ are formed according to \assumptionid{Primary-AM}.
Below is the path diagram for a single pair that assortatively mated on the focal trait with strength $\rho_y$.
The trait's genetic component is the additive sum of two causal variants, which is the smallest example that shows a \emph{cross-variant} covariance; the formulas we state work for any number of variants.

One chain is traced on the diagram to show the rules in use.
It runs from a single allele in one mate to a single allele at a \emph{different} variant in the other --- a covariance that is zero under random mating, and for which there is no common ancestor to appeal to.
Multiplying the coefficients along it gives
$$
    \Cov\!\big[z^{(m)}_{m,1},\, z^{(p)}_{p,2}\big]
    \;=\; \tfrac12 \cdot \beta_1 \cdot \muo \cdot \beta_2 \cdot \tfrac12
    \;=\; \frac{\beta_1\beta_2\,\rho_y}{4\,\VYo} ,
$$
Because there are four such allele pairs behind any one pair of genotypes, they sum back to $\Cov[x_{m,1}, x_{p,2}] = \beta_1\beta_2\rho_y/\VYo$.

\begin{figure}[H]
    \begin{floatcard}
        \definecolor{pmC1F77B4}{HTML}{1F77B4}
\definecolor{pmCBBBBBB}{HTML}{BBBBBB}
\begin{tikzpicture}[
  every node/.style={font=\small},
  pmObserved/.style={draw=black, fill=white, rectangle, inner sep=0.09cm, minimum width=0.42cm, minimum height=0.34cm},
  pmLatent/.style={draw=black, fill=white, ellipse, inner sep=0.16cm, minimum width=0.42cm, minimum height=0.34cm},
  pmDirected/.style={-{Stealth[length=2mm]}, line width=0.7pt},
  pmBidirected/.style={{Stealth[length=2mm]}-{Stealth[length=2mm]}, line width=0.7pt},
  pmCopath/.style={line width=1.6pt},
  pmLabel/.style={midway, fill=white, inner sep=1pt},
]
  \node[pmLatent] (g_m) at (2.500,1.500) {$g_m$};
  \node[pmLatent] (e_m) at (0.000,1.500) {$e_m$};
  \node[pmLatent] (g_p) at (10.500,1.500) {$g_p$};
  \node[pmLatent] (e_p) at (8.000,1.500) {$e_p$};
  \node[pmLatent] (z_mat_m1) at (0.000,4.600) {$z^{(m)}_{m,1}$};
  \node[pmLatent] (z_mat_m2) at (3.400,4.600) {$z^{(m)}_{m,2}$};
  \node[pmLatent] (z_pat_m1) at (1.600,4.600) {$z^{(p)}_{m,1}$};
  \node[pmLatent] (z_pat_m2) at (5.000,4.600) {$z^{(p)}_{m,2}$};
  \node[pmLatent] (z_mat_p1) at (8.000,4.600) {$z^{(m)}_{p,1}$};
  \node[pmLatent] (z_mat_p2) at (11.400,4.600) {$z^{(m)}_{p,2}$};
  \node[pmLatent] (z_pat_p1) at (9.600,4.600) {$z^{(p)}_{p,1}$};
  \node[pmLatent] (z_pat_p2) at (13.000,4.600) {$z^{(p)}_{p,2}$};
  \node[pmObserved] (x_m1) at (0.800,3.000) {$x_{m,1}$};
  \node[pmObserved] (x_m2) at (4.200,3.000) {$x_{m,2}$};
  \node[pmObserved] (x_p1) at (8.800,3.000) {$x_{p,1}$};
  \node[pmObserved] (x_p2) at (12.200,3.000) {$x_{p,2}$};
  \node[pmObserved] (y_m) at (2.500,0.000) {$y_m$};
  \node[pmObserved] (y_p) at (10.500,0.000) {$y_p$};
  \draw[pmDirected, draw=pmC1F77B4, line width=1.82pt] (z_mat_m1) -- node[pmLabel, text=pmC1F77B4] {$1$} (x_m1);
  \draw[pmDirected, draw=pmCBBBBBB] (z_pat_m1) -- node[pmLabel, text=pmCBBBBBB] {$1$} (x_m1);
  \draw[pmDirected, draw=pmCBBBBBB] (z_mat_m2) -- node[pmLabel, text=pmCBBBBBB] {$1$} (x_m2);
  \draw[pmDirected, draw=pmCBBBBBB] (z_pat_m2) -- node[pmLabel, text=pmCBBBBBB] {$1$} (x_m2);
  \draw[pmDirected, draw=pmCBBBBBB] (z_mat_p1) -- node[pmLabel, text=pmCBBBBBB] {$1$} (x_p1);
  \draw[pmDirected, draw=pmCBBBBBB] (z_pat_p1) -- node[pmLabel, text=pmCBBBBBB] {$1$} (x_p1);
  \draw[pmDirected, draw=pmCBBBBBB] (z_mat_p2) -- node[pmLabel, text=pmCBBBBBB] {$1$} (x_p2);
  \draw[pmDirected, draw=pmC1F77B4, line width=1.82pt] (z_pat_p2) -- node[pmLabel, text=pmC1F77B4] {$1$} (x_p2);
  \draw[pmDirected, draw=pmC1F77B4, line width=1.82pt] (x_m1) -- node[pmLabel, text=pmC1F77B4] {$\beta_{1}$} (g_m);
  \draw[pmDirected, draw=pmCBBBBBB] (x_m2) -- node[pmLabel, text=pmCBBBBBB] {$\beta_{2}$} (g_m);
  \draw[pmDirected, draw=pmCBBBBBB] (x_p1) -- node[pmLabel, text=pmCBBBBBB] {$\beta_{1}$} (g_p);
  \draw[pmDirected, draw=pmC1F77B4, line width=1.82pt] (x_p2) -- node[pmLabel, text=pmC1F77B4] {$\beta_{2}$} (g_p);
  \draw[pmDirected, draw=pmC1F77B4, line width=1.82pt] (g_m) -- node[pmLabel, text=pmC1F77B4] {$1$} (y_m);
  \draw[pmDirected, draw=pmCBBBBBB] (e_m) -- node[pmLabel, text=pmCBBBBBB] {$1$} (y_m);
  \draw[pmDirected, draw=pmC1F77B4, line width=1.82pt] (g_p) -- node[pmLabel, text=pmC1F77B4] {$1$} (y_p);
  \draw[pmDirected, draw=pmCBBBBBB] (e_p) -- node[pmLabel, text=pmCBBBBBB] {$1$} (y_p);
  \draw[pmBidirected, draw=pmC1F77B4, line width=1.82pt] (z_mat_m1) to[loop above, looseness=6, in=120, out=60] node[pmLabel, text=pmC1F77B4] {$\frac{1}{2}$} (z_mat_m1);
  \draw[pmBidirected, draw=pmCBBBBBB] (z_mat_m2) to[loop above, looseness=6, in=120, out=60] node[pmLabel, text=pmCBBBBBB] {$\frac{1}{2}$} (z_mat_m2);
  \draw[pmBidirected, draw=pmCBBBBBB] (z_pat_m1) to[loop above, looseness=6, in=120, out=60] node[pmLabel, text=pmCBBBBBB] {$\frac{1}{2}$} (z_pat_m1);
  \draw[pmBidirected, draw=pmCBBBBBB] (z_pat_m2) to[loop above, looseness=6, in=120, out=60] node[pmLabel, text=pmCBBBBBB] {$\frac{1}{2}$} (z_pat_m2);
  \draw[pmBidirected, draw=pmCBBBBBB] (z_mat_p1) to[loop above, looseness=6, in=120, out=60] node[pmLabel, text=pmCBBBBBB] {$\frac{1}{2}$} (z_mat_p1);
  \draw[pmBidirected, draw=pmCBBBBBB] (z_mat_p2) to[loop above, looseness=6, in=120, out=60] node[pmLabel, text=pmCBBBBBB] {$\frac{1}{2}$} (z_mat_p2);
  \draw[pmBidirected, draw=pmCBBBBBB] (z_pat_p1) to[loop above, looseness=6, in=120, out=60] node[pmLabel, text=pmCBBBBBB] {$\frac{1}{2}$} (z_pat_p1);
  \draw[pmBidirected, draw=pmC1F77B4, line width=1.82pt] (z_pat_p2) to[loop above, looseness=6, in=120, out=60] node[pmLabel, text=pmC1F77B4] {$\frac{1}{2}$} (z_pat_p2);
  \draw[pmBidirected, draw=pmCBBBBBB] (e_m) to[loop above, looseness=6, in=120, out=60] node[pmLabel, text=pmCBBBBBB] {$V_{E}$} (e_m);
  \draw[pmBidirected, draw=pmCBBBBBB] (e_p) to[loop above, looseness=6, in=120, out=60] node[pmLabel, text=pmCBBBBBB] {$V_{E}$} (e_p);
  \draw[pmCopath, draw=pmC1F77B4, line width=4.16pt] (y_m) -- node[pmLabel, text=pmC1F77B4] {$\left[\rho_{y}\right]$} (y_p);
  \node[align=center] at (6.500,-1.350) {$z^{(m)}_{m,1} \leftrightarrow z^{(m)}_{m,1} \rightarrow x_{m,1} \rightarrow g_m \rightarrow y_m \;\text{---}\; y_p \leftarrow g_p \leftarrow x_{p,2} \leftarrow z^{(p)}_{p,2} \leftrightarrow z^{(p)}_{p,2}$\\$\operatorname{Cov}\left[z^{(m)}_{m,1}, z^{(p)}_{p,2}\right] = \frac{1}{2} \cdot 1 \cdot \beta_{1} \cdot 1 \cdot \frac{\rho_{y}}{\VPo} \cdot 1 \cdot \beta_{2} \cdot 1 \cdot \frac{1}{2} = \frac{\beta_{1} \beta_{2} \rho_{y}}{4 \VPo}$};
\end{tikzpicture}

        \caption{%
            A mated pair in the base population, assorting on the phenotype with strength $\rho_y$, with one chain traced on it and the product of its coefficients written underneath.
            The model is built from the haploid alleles: $z^{(m)}_{i,k}$ and $z^{(p)}_{i,k}$ are the maternally and paternally inherited alleles of individual $i$ at variant $k$, each with variance $\tfrac12$, and the genotype $x_{i,k} = z^{(m)}_{i,k} + z^{(p)}_{i,k}$.
            Then $g_i = \beta_1 x_{i,1} + \beta_2 x_{i,2}$ and $y_i = g_i + e_i$.
            Genotypes and phenotypes are observed (rectangles); alleles, genetic and environmental values are latent (ellipses) --- an allele's parental origin is not read off an unphased genotype.
            The single thick line is the co-path, labelled with the correlation it induces.
            The edges the chain does not use are faded rather than omitted, so the figure still states the whole model and any other covariance can be traced off the same picture.%
        }
        \label{fig:pair-2v}
    \end{floatcard}
\end{figure}

Doing the same for every pair of variables in the diagram gives the two tables below.
Within an individual (Table~\ref{tab:pair-2v-within}) nothing has changed from the randomly mating case:
these are base-population individuals, so their genotypes are still uncorrelated with each other and with their own environment.

\begin{table}[H]
    \begin{floatcard}
    \begin{tabular}{l|ccccc}
                  & $x_{i,1}$      & $x_{i,2}$      & $g_i$     & $e_i$   & $y_i$     \\
        \hline
        $x_{i,1}$ & $1$            & $0$            & $\beta_1$ & $0$     & $\beta_1$ \\
        $x_{i,2}$ & $0$            & $1$            & $\beta_2$ & $0$     & $\beta_2$ \\
        $g_i$     & $\beta_1$      & $\beta_2$      & $\VAo$    & $0$     & $\VAo$    \\
        $e_i$     & $0$            & $0$            & $0$       & $V_E$   & $V_E$     \\
        $y_i$     & $\beta_1$      & $\beta_2$      & $\VAo$    & $V_E$   & $\VYo$    \\
    \end{tabular}
    \caption{%
        Covariances within one individual $i \in \{m,p\}$, identical for both members of the pair.
        Here $\VAo = \beta_1^2 + \beta_2^2$ and $\VYo = \VAo + V_E$.
        The last column --- each variable's covariance with its own phenotype --- is the only part of this table needed for computing the covariance between variables across mates.%
    }
    \label{tab:pair-2v-within}
    \end{floatcard}
\end{table}

The co-path induces covariance between the factors pertaining to mates such that:
\begin{equation}\label{eq:leg-rule}
    \boxed{\ \Cov[\,a_m,\, b_p\,] \;=\; \Cov[\,a_m,\, y_m\,] \cdot \muo \cdot \Cov[\,y_p,\, b_p\,] \ }
\end{equation}
for any variables $a$ and $b$ belonging to the maternal and paternal member respectively.
The entire cross-mate block is therefore the outer product of a single vector --- the last column of Table~\ref{tab:pair-2v-within} --- scaled by $\muo = \rho_y/\VYo$.

\begin{table}[H]
    \begin{floatcard}
    \begin{tabular}{l|ccccc}
        $\times \; \muo$ & $x_{p,1}$       & $x_{p,2}$         & $g_p$            & $e_p$          & $y_p$            \\
        \hline
        $x_{m,1}$ & $\beta_1^2$        & $\beta_1\beta_2$   & $\beta_1\VAo$     & $\beta_1 V_E$   & $\beta_1\VYo$     \\
        $x_{m,2}$ & $\beta_1\beta_2$   & $\beta_2^2$        & $\beta_2\VAo$     & $\beta_2 V_E$   & $\beta_2\VYo$     \\
        $g_m$     & $\beta_1\VAo$      & $\beta_2\VAo$      & $(\VAo)^2$        & $\VAo V_E$      & $\VAo\VYo$        \\
        $e_m$     & $\beta_1 V_E$      & $\beta_2 V_E$      & $\VAo V_E$        & $V_E^2$         & $V_E\VYo$         \\
        $y_m$     & $\beta_1\VYo$      & $\beta_2\VYo$      & $\VAo\VYo$        & $V_E\VYo$       & $(\VYo)^2$        \\
    \end{tabular}
    \caption{%
        Covariances between the two mates. \textbf{Every entry carries a factor $\muo = \rho_y/\VYo$}, shown once in the corner rather than repeated in twenty-five cells.
        The block is symmetric, and is the outer product of $(\beta_1, \beta_2, \VAo, V_E, \VYo)$ with itself --- a direct consequence of \eqref{eq:leg-rule}.
        Setting $\rho_y = 0$ returns every entry to zero, recovering the randomly mating case of Section~1.%
    }
    \label{tab:pair-2v-between}
    \end{floatcard}
\end{table}

Three key results from Table~\ref{tab:pair-2v-between} are highlighted here.
We also generalize from two variants to $M$ variants, where $g_i = \sum_{k=1}^{M} \beta_k x_{ik}$.

\begin{itemize}
    \item $\Cov[\,x_{mk},\, x_{pl}\,] = \beta_k\beta_l\,\rho_y/\VYo \neq 0$, for \emph{every} pair of variants $k$ and $l$.
    Alleles at different loci in different individuals are correlated, signed by the product of their effects and breaking the \assumptionid{Non-IBD-indep} assumption.
    Note that in the base population, the within-individual covariance $\Cov[\,x_{mk},\, x_{ml}\,]$ is still exactly zero for $k \neq l$; the disequilibrium exists only \emph{between} the mates.
    \item $\Cov[\,e_m,\, g_p\,] = \VAo V_E \,\rho_y/\VYo \neq 0$.
    One mate's environment is correlated with the other's genetic value, which is the break of \assumptionid{Env-indep} and \assumptionid{No-IGE} flagged at the start of this section.
    Despite one mate's genotype not causally affecting the other's environment, the two still become correlated due to the phenotypic matching.
    \item $\rhogo \,=\, \Corr[\,g_m,\, g_p\,] \,=\, \frac{(\VAo)^2 \rho_y/\VYo}{\VAo} \,=\, \rho_y\,\frac{\VAo}{\VYo} \,=\, \rho_y\,\hsqo . $
    Mates' genetic values are correlated according to the strength of assortative mating and the heritability of the sorting trait.
    Equivalently, this is the sum of the first result over all $M^2$ ordered pairs of variants, weighted by their effects: $\sum_{k}\sum_{l}\beta_k\beta_l \cdot \beta_k\beta_l\rho_y/\VYo = (\sum_k \beta_k^2)^2\rho_y/\VYo$.
\end{itemize}

\subsubsection{Consequences of assortative mating on first generation of offspring}
Next, we will extend our path diagram to contain two offspring of an assortatively mated pair.
For each variant, all offspring inherit one allele from each parent, but which allele they inherit from that parent (the maternal vs paternal allele) is uniformly random.
For the purposes of tracking covariances, we model this by drawing an edge from each parental allele to their respective inherited allele in the offspring and adding an exogeneous source of variance.
The directed edges have a coefficient of $\tfrac12$ to reflect that the expected inherited allele in the offspring is the average of the two alleles of the parent being inherited from.

Write $o$ for an offspring of the mated pair, and let $B^{(m)}_{ok} \in \{0,1\}$ indicate which of the mother's two alleles was transmitted at variant $k$.
We arbitrarily label the maternal allele $z^{(m)}_{mk}$ and the paternal allele $z^{(p)}_{mk}$ such that if $B^{(m)}_{ok} = 1$, then the offspring inherits the maternal allele, and if $B^{(m)}_{ok} = 0$ then the offspring inherits the paternal allele.
\begin{equation}\label{eq:transmit}
    z^{(m)}_{ok} \;=\; B^{(m)}_{ok}\, z^{(m)}_{mk} \;+\; \big(1 - B^{(m)}_{ok}\big)\, z^{(p)}_{mk} ,
    \qquad B^{(m)}_{ok} \sim \mathrm{Bernoulli}(\tfrac12), \quad B^{(m)}_{ok} \perp z_{mk} .
\end{equation}
A path diagram represents the \emph{linear projection} of the random variable: the best prediction of the offspring's allele from the mother's two, which is their average,
\begin{equation}\label{eq:transmit-mean}
    \E\!\big[\,z^{(m)}_{ok} \,\big|\, z^{(m)}_{mk}, z^{(p)}_{mk}\,\big]
    \;=\; \tfrac12 z^{(m)}_{mk} + \tfrac12 z^{(p)}_{mk}
    \;=\; \tfrac12 x_{mk} .
\end{equation}

The exogenous variance is what the projection leaves behind.
Subtracting \eqref{eq:transmit-mean} from \eqref{eq:transmit} gives the \emph{segregation deviation}
\begin{equation}\label{eq:seg-def}
    z^{\perp(m)}_{ok} \;:=\; z^{(m)}_{ok} - \tfrac12 x_{mk} \;=\; \big(B^{(m)}_{ok} - \tfrac12\big)\big(z^{(m)}_{mk} - z^{(p)}_{mk}\big) ,
\end{equation}
which is mean-zero and uncorrelated with both of the mother's alleles --- a disturbance in the sense of Section~2.2.1, and drawn as one.
The offspring's paternally inherited allele $z^{(p)}_{ok}$ is built the same way from the father's two alleles, with its own independent $B^{(p)}_{ok}$ and residual $z^{\perp(p)}_{ok}$.

\subsubsection*{Covariance between alleles}
We now define the covariances that will be repeatedly used in deriving the consequences of assortartive mating.
For an individual $i$ of generation $t$ and variants $k$ and $l$ where $u,v \in \{m,p\}$,
\begin{equation}\label{eq:as-def}
    \boxed{\
    a^{(t)}_{kl} \;:=\; \Cov\!\big[z^{(u)}_{ik},\, z^{(v)}_{il}\big]\ \ (u \neq v),
    \qquad
    s^{(t)}_{kl} \;:=\; \Cov\!\big[z^{(u)}_{ik},\, z^{(u)}_{il}\big] \ }
\end{equation}
In words, $a$ denotes the covariance between alleles \emph{across} gametes, while $s$ denotes the covariance between alleles on the \emph{same} gamete.
When $k=l$, $s^{(t)}_{kk}$ is the variance of an individual allele, which is fixed for all generations due to how we standardized variables according to \assumptionid{Constant-freq}:
\begin{equation}\label{eq:skk}
    s^{(t)}_{kk} \;\equiv\; \tfrac12 \qquad \text{for all } t ,
\end{equation}
We can also write the covariance between two genotypes in the same individual as the sum of four allelic covariances:
\begin{equation}\label{eq:x-from-as}
    \Cov[\,x_{ik},\, x_{il}\,]
    \;=\; \sum_{u,v \in \{m,p\}} \Cov\!\big[z^{(u)}_{ik},\, z^{(v)}_{il}\big]
    \;=\; 2\big(s^{(t)}_{kl} + a^{(t)}_{kl}\big) ,
\end{equation}
In the case of $k=l$, it becomes the variance of genotypes, since $s^{(t)}_{kk} \equiv \tfrac12$:
\begin{equation}\label{eq:varx-a}
\begin{gathered}
    \Var[\,x_{ik}\,] \;=\; 2\big(s^{(t)}_{kk} + a^{(t)}_{kk}\big) \\[4pt]
    \boxed{\ \Var[\,x_{ik}\,] \;=\; 1 + 2\,a^{(t)}_{kk}\ }
\end{gathered}
\end{equation}

The third quantity we need is the covariance between a single allele and the phenotype of the individual carrying it.
Expanding the phenotype into its causes and then each genotype into its two alleles,
\begin{equation}\label{eq:c-def}
    \boxed{\ c^{(t)}_k \;:=\; \Cov\!\big[\,z^{(u)}_{ik},\, y_i\,\big]
    \;=\; \sum_{l=1}^{M} \beta_l\,\Cov\!\big[\,z^{(u)}_{ik},\, x_{il}\,\big]
    \;=\; \sum_{l=1}^{M} \beta_l\Big(s^{(t)}_{kl} + a^{(t)}_{kl}\Big) \ }
\end{equation}
The first step drops $e_i$, which is uncorrelated with every allele of the same individual by \assumptionid{GE-indep}.
The second splits each genotype $x_{il}$ into the allele sitting on the same gamete as $z^{(u)}_{ik}$, contributing $s^{(t)}_{kl}$, and the one on the opposite gamete, contributing $a^{(t)}_{kl}$.
A single allele therefore reaches the phenotype through \emph{every} variant the individual carries, not only through its own.

Summing over alleles in the same way gives the additive variance of the generation:
\begin{equation}\label{eq:VA-from-as}
\begin{gathered}
\begin{aligned}
    \VAt \;=\; \Var\Big[\textstyle\sum_k \beta_k x_{ik}\Big]
    \;&=\; \sum_{k}\sum_{l} \beta_k\beta_l \Cov[\,x_{ik},\, x_{il}\,] \\[2pt]
    \;&=\; 2\sum_{k}\sum_{l} \beta_k\beta_l\big(s^{(t)}_{kl} + a^{(t)}_{kl}\big)
\end{aligned} \\[6pt]
    \boxed{\ \VAt \;=\; 2\sum_{k} \beta_k\, c^{(t)}_k \ }
\end{gathered}
\end{equation}
with $\VYt = \VAt + V_E$, since assortment on the phenotype does not change the environmental variance.
Similarly, we define generation-specific heritability as follows:
\begin{equation}\label{eq:hsq-t}
    \boxed{\ \hsqt \;:=\; \frac{\VAt}{\VYt} \;=\; \frac{\VAt}{\VAt + V_E}
    \;=\; \frac{2\sum_k \beta_k c^{(t)}_k}{2\sum_k \beta_k c^{(t)}_k + V_E} \ }
\end{equation}

We'll also confirm here that $\rho_g^{(t)} = \frac{ \rho_y V_A^{(t)} }{V_Y^{(t)} }$.
A single co-path joins the two mates, so \eqref{eq:leg-rule} applies with $\mut$ in place of $\muo$, and a genetic value's covariance with its own carrier's phenotype is $\Cov[\,g_i,\, g_i + e_i\,] = \VAt$ by \assumptionid{GE-indep}.
Both mates have genetic variance $\VAt$, so the denominator of the correlation is $\VAt$ as well:
\begin{equation}\label{eq:rhog-t}
\begin{gathered}
    \rhogt \;=\; \frac{\Cov[\,g_m,\, g_p\,]}{\sqrt{\Var[g_m]\,\Var[g_p]}}
    \;=\; \frac{\overbrace{\Cov[\,g_m,\, y_m\,]}^{\VAt} \cdot\; \mut \;\cdot \overbrace{\Cov[\,y_p,\, g_p\,]}^{\VAt}}{\VAt}
    \;=\; \mut\,\VAt \\[4pt]
    \boxed{\ \rhogt \;=\; \frac{\rho_y\,\VAt}{\VYt} \;=\; \rho_y\,\hsqt \ }
\end{gathered}
\end{equation}
The genetic correlation between mates is the phenotypic one scaled by that generation's heritability, exactly as it was in the base population.

We can also now write the variance of the segregation deviation.
The indicator $B^{(m)}_{ok}$ takes values in $\{0,1\}$, so $\big(B^{(m)}_{ok} - \tfrac12\big)^2 = \tfrac14$ with probability one, and it is independent of the alleles it selects between.
Squaring \eqref{eq:seg-def} and taking expectations therefore factors cleanly:
\begin{equation}\label{eq:seg-var-deriv}
    \Var\!\big[z^{\perp(m)}_{ok}\big]
    \;=\; \E\!\Big[\big(B^{(m)}_{ok} - \tfrac12\big)^2\Big]\; \Var\!\big[z^{(m)}_{mk} - z^{(p)}_{mk}\big]
    \;=\; \tfrac14\Big(\underbrace{\tfrac12}_{s^{(t)}_{kk}} + \underbrace{\tfrac12}_{s^{(t)}_{kk}} - 2\,\underbrace{a^{(t)}_{kk}}_{\Cov[z^{(m)}_{mk},\, z^{(p)}_{mk}]}\Big) ,
\end{equation}
where the parent is of generation $t$. Collecting terms,
\begin{equation}\label{eq:seg-var-a}
    \boxed{\ \Var\!\big[z^{\perp(m)}_{ok}\big] \;=\; \tfrac14 - \tfrac12\,a^{(t)}_{kk} \ }
\end{equation}
So a parent whose two alleles are more alike contributes \emph{less} segregation variance: it matters less which of the two was transmitted.

\subsubsection*{First generation of offspring}

In a generation-0 individual, their own alleles are uncorrelated with each other.
At the same variant, the two alleles are non-IBD by \assumptionid{No-inbreeding} and hence independent by \assumptionid{Non-IBD-indep}; at different variants, the base population is in linkage equilibrium, which is the within-individual zero already noted below Table~\ref{tab:pair-2v-between}:
\begin{equation}\label{eq:as-gen0}
    \boxed{\ a^{(0)}_{kl} \;=\; 0 \ \ \text{for all } k, l ,
    \qquad
    s^{(0)}_{kl} \;=\; 0 \ \ (k \neq l),
    \qquad
    c^{(0)}_k \;=\; \tfrac12\beta_k \ }
\end{equation}
The third follows from the first two: with every off-diagonal term of \eqref{eq:c-def} equal to zero, only $l=k$ survives, leaving $\beta_k s^{(0)}_{kk} = \tfrac12\beta_k$.
Setting $a^{(0)}_{kk} = 0$ in \eqref{eq:varx-a} recovers $\Var[x_{ik}] = 1$, and in \eqref{eq:seg-var-a} recovers the familiar segregation variance of $\tfrac14$.

In a generation-1 individual, this is no longer the case, which can be seen by the below path diagram of two generation-1 offspring of a mated pair that assortatively mated on the focal trait with strength $\rho_y$, for a trait with two causal variants.

\begin{sidewaysfigure}
    \begin{floatcard}
        \definecolor{pmCB03A2E}{HTML}{B03A2E}
\begin{tikzpicture}[
  every node/.style={font=\small},
  pmObserved/.style={draw=black, fill=white, rectangle, inner sep=0.09cm, minimum width=0.42cm, minimum height=0.34cm},
  pmLatent/.style={draw=black, fill=white, ellipse, inner sep=0.16cm, minimum width=0.42cm, minimum height=0.34cm},
  pmDirected/.style={-{Stealth[length=2mm]}, line width=0.7pt},
  pmBidirected/.style={{Stealth[length=2mm]}-{Stealth[length=2mm]}, line width=0.7pt},
  pmCopath/.style={line width=1.6pt},
  pmLabel/.style={midway, fill=white, inner sep=1pt},
]
  \node[pmLatent] (g_m) at (3.300,11.600) {$g_m$};
  \node[pmLatent] (e_m) at (6.900,11.600) {$e_m$};
  \node[pmLatent] (g_p) at (17.900,11.600) {$g_p$};
  \node[pmLatent] (e_p) at (21.500,11.600) {$e_p$};
  \node[pmLatent] (z_mat_m1) at (0.000,8.400) {$z^{(m)}_{m,1}$};
  \node[pmLatent] (z_mat_m2) at (4.600,8.400) {$z^{(m)}_{m,2}$};
  \node[pmLatent] (z_pat_m1) at (2.000,8.400) {$z^{(p)}_{m,1}$};
  \node[pmLatent] (z_pat_m2) at (6.600,8.400) {$z^{(p)}_{m,2}$};
  \node[pmLatent] (z_mat_p1) at (14.600,8.400) {$z^{(m)}_{p,1}$};
  \node[pmLatent] (z_mat_p2) at (19.200,8.400) {$z^{(m)}_{p,2}$};
  \node[pmLatent] (z_pat_p1) at (16.600,8.400) {$z^{(p)}_{p,1}$};
  \node[pmLatent] (z_pat_p2) at (21.200,8.400) {$z^{(p)}_{p,2}$};
  \node[pmLatent] (z_mat_o11) at (0.150,5.000) {$z^{(m)}_{o1,1}$};
  \node[pmLatent] (z_mat_o12) at (4.750,5.000) {$z^{(m)}_{o1,2}$};
  \node[pmLatent] (z_pat_o11) at (14.750,5.000) {$z^{(p)}_{o1,1}$};
  \node[pmLatent] (z_pat_o12) at (19.350,5.000) {$z^{(p)}_{o1,2}$};
  \node[pmLatent] (z_mat_o21) at (1.850,5.000) {$z^{(m)}_{o2,1}$};
  \node[pmLatent] (z_mat_o22) at (6.450,5.000) {$z^{(m)}_{o2,2}$};
  \node[pmLatent] (z_pat_o21) at (16.450,5.000) {$z^{(p)}_{o2,1}$};
  \node[pmLatent] (z_pat_o22) at (21.050,5.000) {$z^{(p)}_{o2,2}$};
  \node[pmLatent] (g_o1) at (6.700,1.100) {$g_{o1}$};
  \node[pmLatent] (g_o2) at (14.500,1.100) {$g_{o2}$};
  \node[pmLatent] (e_o1) at (3.300,1.100) {$e_{o1}$};
  \node[pmLatent] (e_o2) at (11.100,1.100) {$e_{o2}$};
  \node[pmObserved] (x_m1) at (1.000,10.000) {$x_{m,1}$};
  \node[pmObserved] (x_m2) at (5.600,10.000) {$x_{m,2}$};
  \node[pmObserved] (x_p1) at (15.600,10.000) {$x_{p,1}$};
  \node[pmObserved] (x_p2) at (20.200,10.000) {$x_{p,2}$};
  \node[pmObserved] (y_m) at (3.300,13.200) {$y_m$};
  \node[pmObserved] (y_p) at (17.900,13.200) {$y_p$};
  \node[pmObserved] (x_o11) at (5.200,2.600) {$x_{o1,1}$};
  \node[pmObserved] (x_o12) at (8.200,2.600) {$x_{o1,2}$};
  \node[pmObserved] (y_o1) at (6.700,-0.300) {$y_{o1}$};
  \node[pmObserved] (x_o21) at (13.000,2.600) {$x_{o2,1}$};
  \node[pmObserved] (x_o22) at (16.000,2.600) {$x_{o2,2}$};
  \node[pmObserved] (y_o2) at (14.500,-0.300) {$y_{o2}$};
  \draw[pmDirected, draw=black] (z_mat_m1) -- (x_m1);
  \draw[pmDirected, draw=black] (z_pat_m1) -- (x_m1);
  \draw[pmDirected, draw=black] (z_mat_m2) -- (x_m2);
  \draw[pmDirected, draw=black] (z_pat_m2) -- (x_m2);
  \draw[pmDirected, draw=black] (z_mat_p1) -- (x_p1);
  \draw[pmDirected, draw=black] (z_pat_p1) -- (x_p1);
  \draw[pmDirected, draw=black] (z_mat_p2) -- (x_p2);
  \draw[pmDirected, draw=black] (z_pat_p2) -- (x_p2);
  \draw[pmDirected, draw=black] (x_m1) -- node[pmLabel] {$\beta_{1}$} (g_m);
  \draw[pmDirected, draw=black] (x_m2) -- node[pmLabel] {$\beta_{2}$} (g_m);
  \draw[pmDirected, draw=black] (x_p1) -- node[pmLabel] {$\beta_{1}$} (g_p);
  \draw[pmDirected, draw=black] (x_p2) -- node[pmLabel] {$\beta_{2}$} (g_p);
  \draw[pmDirected, draw=black] (g_m) -- (y_m);
  \draw[pmDirected, draw=black] (e_m) -- (y_m);
  \draw[pmDirected, draw=black] (g_p) -- (y_p);
  \draw[pmDirected, draw=black] (e_p) -- (y_p);
  \draw[pmDirected, draw=black] (z_mat_m1) -- node[pmLabel] {$\frac{1}{2}$} (z_mat_o11);
  \draw[pmDirected, draw=black] (z_pat_m1) -- node[pmLabel] {$\frac{1}{2}$} (z_mat_o11);
  \draw[pmDirected, draw=black] (z_mat_m2) -- node[pmLabel] {$\frac{1}{2}$} (z_mat_o12);
  \draw[pmDirected, draw=black] (z_pat_m2) -- node[pmLabel] {$\frac{1}{2}$} (z_mat_o12);
  \draw[pmDirected, draw=black] (z_mat_p1) -- node[pmLabel] {$\frac{1}{2}$} (z_pat_o11);
  \draw[pmDirected, draw=black] (z_pat_p1) -- node[pmLabel] {$\frac{1}{2}$} (z_pat_o11);
  \draw[pmDirected, draw=black] (z_mat_p2) -- node[pmLabel] {$\frac{1}{2}$} (z_pat_o12);
  \draw[pmDirected, draw=black] (z_pat_p2) -- node[pmLabel] {$\frac{1}{2}$} (z_pat_o12);
  \draw[pmDirected, draw=black] (z_mat_o11) -- (x_o11);
  \draw[pmDirected, draw=black] (z_pat_o11) -- (x_o11);
  \draw[pmDirected, draw=black] (z_mat_o12) -- (x_o12);
  \draw[pmDirected, draw=black] (z_pat_o12) to[bend left=10] (x_o12);
  \draw[pmDirected, draw=black] (x_o11) -- node[pmLabel] {$\beta_{1}$} (g_o1);
  \draw[pmDirected, draw=black] (x_o12) -- node[pmLabel] {$\beta_{2}$} (g_o1);
  \draw[pmDirected, draw=black] (g_o1) -- (y_o1);
  \draw[pmDirected, draw=black] (e_o1) -- (y_o1);
  \draw[pmDirected, draw=black] (z_mat_m1) -- node[pmLabel, pos=0.380, xshift=0.351cm, yshift=0.191cm] {$\frac{1}{2}$} (z_mat_o21);
  \draw[pmDirected, draw=black] (z_pat_m1) -- node[pmLabel] {$\frac{1}{2}$} (z_mat_o21);
  \draw[pmDirected, draw=black] (z_mat_m2) -- node[pmLabel, pos=0.380, xshift=0.351cm, yshift=0.191cm] {$\frac{1}{2}$} (z_mat_o22);
  \draw[pmDirected, draw=black] (z_pat_m2) -- node[pmLabel] {$\frac{1}{2}$} (z_mat_o22);
  \draw[pmDirected, draw=black] (z_mat_p1) -- node[pmLabel, pos=0.380, xshift=0.351cm, yshift=0.191cm] {$\frac{1}{2}$} (z_pat_o21);
  \draw[pmDirected, draw=black] (z_pat_p1) -- node[pmLabel] {$\frac{1}{2}$} (z_pat_o21);
  \draw[pmDirected, draw=black] (z_mat_p2) -- node[pmLabel, pos=0.380, xshift=0.351cm, yshift=0.191cm] {$\frac{1}{2}$} (z_pat_o22);
  \draw[pmDirected, draw=black] (z_pat_p2) -- node[pmLabel] {$\frac{1}{2}$} (z_pat_o22);
  \draw[pmDirected, draw=black] (z_mat_o21) to[bend right=10] (x_o21);
  \draw[pmDirected, draw=black] (z_pat_o21) -- (x_o21);
  \draw[pmDirected, draw=black] (z_mat_o22) -- (x_o22);
  \draw[pmDirected, draw=black] (z_pat_o22) -- (x_o22);
  \draw[pmDirected, draw=black] (x_o21) -- node[pmLabel] {$\beta_{1}$} (g_o2);
  \draw[pmDirected, draw=black] (x_o22) -- node[pmLabel] {$\beta_{2}$} (g_o2);
  \draw[pmDirected, draw=black] (g_o2) -- (y_o2);
  \draw[pmDirected, draw=black] (e_o2) -- (y_o2);
  \draw[pmBidirected, draw=black] (z_mat_m1) to[loop above, looseness=6, in=120, out=60] node[pmLabel] {$\frac{1}{2}$} (z_mat_m1);
  \draw[pmBidirected, draw=black] (z_mat_m2) to[loop above, looseness=6, in=120, out=60] node[pmLabel] {$\frac{1}{2}$} (z_mat_m2);
  \draw[pmBidirected, draw=black] (z_pat_m1) to[loop above, looseness=6, in=120, out=60] node[pmLabel] {$\frac{1}{2}$} (z_pat_m1);
  \draw[pmBidirected, draw=black] (z_pat_m2) to[loop above, looseness=6, in=120, out=60] node[pmLabel] {$\frac{1}{2}$} (z_pat_m2);
  \draw[pmBidirected, draw=black] (z_mat_p1) to[loop above, looseness=6, in=120, out=60] node[pmLabel] {$\frac{1}{2}$} (z_mat_p1);
  \draw[pmBidirected, draw=black] (z_mat_p2) to[loop above, looseness=6, in=120, out=60] node[pmLabel] {$\frac{1}{2}$} (z_mat_p2);
  \draw[pmBidirected, draw=black] (z_pat_p1) to[loop above, looseness=6, in=120, out=60] node[pmLabel] {$\frac{1}{2}$} (z_pat_p1);
  \draw[pmBidirected, draw=black] (z_pat_p2) to[loop above, looseness=6, in=120, out=60] node[pmLabel] {$\frac{1}{2}$} (z_pat_p2);
  \draw[pmBidirected, draw=black] (e_m) to[loop above, looseness=6, in=120, out=60] node[pmLabel] {$V_{E}$} (e_m);
  \draw[pmBidirected, draw=black] (e_p) to[loop above, looseness=6, in=120, out=60] node[pmLabel] {$V_{E}$} (e_p);
  \draw[pmBidirected, draw=black] (z_mat_o11) to[loop above, looseness=6, in=120, out=60] node[pmLabel] {$\frac{1}{4}$} (z_mat_o11);
  \draw[pmBidirected, draw=black] (z_mat_o12) to[loop above, looseness=6, in=120, out=60] node[pmLabel] {$\frac{1}{4}$} (z_mat_o12);
  \draw[pmBidirected, draw=black] (z_pat_o11) to[loop above, looseness=6, in=120, out=60] node[pmLabel] {$\frac{1}{4}$} (z_pat_o11);
  \draw[pmBidirected, draw=black] (z_pat_o12) to[loop above, looseness=6, in=120, out=60] node[pmLabel] {$\frac{1}{4}$} (z_pat_o12);
  \draw[pmBidirected, draw=black] (e_o1) to[loop above, looseness=6, in=120, out=60] node[pmLabel] {$V_{E}$} (e_o1);
  \draw[pmBidirected, draw=black] (z_mat_o21) to[loop above, looseness=6, in=120, out=60] node[pmLabel] {$\frac{1}{4}$} (z_mat_o21);
  \draw[pmBidirected, draw=black] (z_mat_o22) to[loop above, looseness=6, in=120, out=60] node[pmLabel] {$\frac{1}{4}$} (z_mat_o22);
  \draw[pmBidirected, draw=black] (z_pat_o21) to[loop above, looseness=6, in=120, out=60] node[pmLabel] {$\frac{1}{4}$} (z_pat_o21);
  \draw[pmBidirected, draw=black] (z_pat_o22) to[loop above, looseness=6, in=120, out=60] node[pmLabel] {$\frac{1}{4}$} (z_pat_o22);
  \draw[pmBidirected, draw=black] (e_o2) to[loop above, looseness=6, in=120, out=60] node[pmLabel] {$V_{E}$} (e_o2);
  \draw[pmCopath, draw=pmCB03A2E] (y_m) -- node[pmLabel] {$\left[\rho_{y}\right]$} (y_p);
\end{tikzpicture}

        \caption{%
            One mated pair from the base population and two of their offspring, at the allele level, for a trait with two causal variants.
            Each inherited allele is fed by \emph{both} alleles of the parent it came from, at $\tfrac12$ each, per \eqref{eq:transmit-mean}.
            The self-loop of $\tfrac14$ on each offspring allele is its segregation variance, \eqref{eq:seg-var-a} evaluated at $a^{(0)}_{kk} = 0$.%
        }
        \label{fig:pair-offspring-2v}
    \end{floatcard}
\end{sidewaysfigure}

We can use our path tracing rules to compute $a^{(1)}_{kl}$.
Exactly four chains connect an allele of the offspring to an allele on its \emph{other} gamete, at variants $k$ and $l$.
Each traces back from $z^{(m)}_{ok}$ to \emph{one} of the mother's two alleles at variant $k$, turns around on that allele's variance, runs up through her genotype and phenotype, crosses the co-path, and comes back down the father's side to \emph{one} of his two alleles at variant $l$ and into $z^{(p)}_{ol}$.
Two choices on each side give $2\times2 = 4$ chains identical in value, one for each $u, v \in \{m,p\}$:
$$
\begin{aligned}
    &z^{(m)}_{ok} \leftarrow z^{(u)}_{mk} \leftrightarrow z^{(u)}_{mk} \rightarrow x_{mk} \rightarrow g_m \rightarrow y_m \\
    \;\text{---}\;\;
    &y_p \leftarrow g_p \leftarrow x_{pl} \leftarrow z^{(v)}_{pl} \leftrightarrow z^{(v)}_{pl} \rightarrow z^{(p)}_{ol} ,
\end{aligned}
$$
Multiplying each coefficient along one chain returns:
$$
    \underbrace{\tfrac12}_{\text{transmission}} \cdot \underbrace{\tfrac12}_{\Var[z]} \cdot \beta_k \cdot \muo \cdot \beta_l \cdot \underbrace{\tfrac12}_{\Var[z]} \cdot \underbrace{\tfrac12}_{\text{transmission}}
    \;=\; \frac{\beta_k\beta_l\,\muo}{16} .
$$
Summing across the four chains and expanding $\muo$,
\begin{equation}\label{eq:a-gen1}
     \boxed{a^{(1)}_{kl} \;=\; \frac{\beta_k\beta_l\,\rho_y}{4\,\VYo} \ }
\end{equation}
The same-gamete companion is easier: two alleles the offspring received from the \emph{same} parent at different variants have no chain at all connecting them, because that parent's own alleles are uncorrelated across variants due to them being of the base population, so
\begin{equation}\label{eq:s-gen1}
    \boxed{\ s^{(1)}_{kl} \;=\; 0 \quad (k \neq l) \ }
\end{equation}
One generation of assortment therefore produces association \emph{across} an individual's two gametes and none within either of them.

Rearranging \eqref{eq:rhog-t} gives $\rho_y/\VYo = \rhogo/\VAo$, which lets \eqref{eq:a-gen1} be restated in terms of how strongly the parents' genetic values were correlated rather than their phenotypes:
\begin{equation}\label{eq:a-gen1-rhog}
    \boxed{\ a^{(1)}_{kl} \;=\; \frac{\beta_k\beta_l}{4\,\VAo}\;\rhogo \ }
\end{equation}
The association between two of the offspring's alleles is the product of their effect sizes, taken as a share of the additive variance, scaled by the genetic correlation between its parents.

Feeding \eqref{eq:a-gen1} and \eqref{eq:s-gen1} into \eqref{eq:c-def} gives the first-generation allele--phenotype covariance.
Only $l = k$ survives from $s^{(1)}$, while every $l$ contributes through $a^{(1)}$:
\begin{equation}\label{eq:c-gen1}
    \begin{gathered}
        c^{(1)}_k
        \;=\; \underbrace{\tfrac12\beta_k}_{l\,=\,k \text{ term of } s^{(1)}}
        \;+\; \sum_{l} \beta_l\,\frac{\beta_k\beta_l\,\rho_y}{4\,\VYo} 
        \;=\; \tfrac12\beta_k\left(1 + \frac{\rho_y\,\VAo}{2\,\VYo}\right) \\[4pt]
        \boxed{\
        c^{(1)}_k  \;=\; \tfrac12\beta_k\left(1 + \tfrac12\rhogo\right)
        \ }
    \end{gathered}
\end{equation}
using $\sum_l \beta_l^2 = \VAo$.
Every variant's initial covariance with the phenotype is inflated by the same factor.

We use the path tracing rules to compute the covariance between an individual and a full sibling and a parent, shown in Table~\ref{tab:relatives-gen1}.
Every entry there turns out to be its random-mating value plus a term proportional to $\rhogo$.

\begin{table}[H]
    \begin{floatcard}
    \renewcommand{\arraystretch}{1.9}
    \setlength{\tabcolsep}{2pt}
    \footnotesize
    \begin{tabular}{l|ccc}
            & self
            & full siblings
            & parent--offspring \\
        \hline
        genotypes, $k = l$
            & $1 + \adds{\dfrac{\beta_k^2}{2\VAo}\rhogo}$
            & $\tfrac12 + \adds{\dfrac{\beta_k^2}{2\VAo}\rhogo}$
            & $\tfrac12 + \adds{\dfrac{\beta_k^2}{2\VAo}\rhogo}$ \\
        genotypes, $k \neq l$
            & $\adds{\dfrac{\beta_k\beta_l}{2\VAo}\rhogo}$
            & $\adds{\dfrac{\beta_k\beta_l}{2\VAo}\rhogo}$
            & $\adds{\dfrac{\beta_k\beta_l}{2\VAo}\rhogo}$ \\
        genetic values
            & $\VAo\big(1 + \adds{\tfrac12\rhogo}\big)$
            & $\tfrac12\VAo\big(1 + \adds{\rhogo}\big)$
            & $\tfrac12\VAo\big(1 + \adds{\rhogo}\big)$ \\
        phenotypes
            & $\VAo\big(1 + \adds{\tfrac12\rhogo}\big) + V_E$
            & $\tfrac12\VAo\big(1 + \adds{\rhogo}\big)$
            & $\tfrac12\VAo\big(1 + \adds{\rho_y}\big)$ \\
    \end{tabular}
    \caption{%
        Variances and covariances for an individual of the \textbf{first generation produced under assortative mating}, against itself and against each kind of first-degree relative: the ``self'' column is the pair $(o_1, o_1)$, ``full siblings'' is $(o_1, o_2)$, and ``parent--offspring'' is $(m, o_1)$. The first row is a genotype against a genotype at the \emph{same} variant ($k = l$), the second at two \emph{different} variants.
        Every entry is its random-mating value plus \adds{what one generation of assortment added}.
        Writing the added terms through $\rhogo$ rather than $\rho_y/\VYo$ leaves the two relationships differing in a single symbol: siblings carry $\rhogo$ at the phenotypic level where parent and offspring carry $\rho_y$.%
    }
    \label{tab:relatives-gen1}
    \end{floatcard}
\end{table}

Note that covariances for parents vs full-siblings only differ at the phenotypic level due to one extra non-zero term that emerges for direct descendants:
$$
\begin{aligned}
    \Cov[\,y_m, y_{o_1}\,]
    \;&=\; \underbrace{\Cov[\,g_m, g_{o_1}\,]}_{\tfrac12\VAo \,+\, \frac{\rho_y (\VAo)^2}{2\,\VYo}}
    \;+\; \underbrace{\Cov[\,e_m, g_{o_1}\,]}_{\frac{\rho_y\,\VAo V_E}{2\,\VYo}} \\[4pt]
    \;&=\; \tfrac12\VAo \;+\; \frac{\rho_y\,\VAo\big(\VAo + V_E\big)}{2\,\VYo}
    \;=\; \tfrac12\VAo\big(1+\rho_y\big) ,
\end{aligned}
$$
where the last step uses $\VYo = \VAo + V_E$.
A parent's environment is correlated with their child's genetic value --- the child carries half the alleles of the assortatively chosen partner, whose genotype is correlated with this parent's environment.

\subsubsection{Compounding across generations}
If we now produce a second generation of offspring, the effects of another generation of assortative mating compound on the existing effects of past assortative mating.
Namely, unlike before, the two alleles that a parent carries are now correlated according to their effect size.

It would be unwieldy to make a diagram for three generations.
Instead we summarize the base population by its impact on the covariance of generation-1 alleles.

\begin{sidewaysfigure}
    \begin{floatcard}
        \definecolor{pmCB03A2E}{HTML}{B03A2E}
\begin{tikzpicture}[
  every node/.style={font=\small},
  pmObserved/.style={draw=black, fill=white, rectangle, inner sep=0.09cm, minimum width=0.42cm, minimum height=0.34cm},
  pmLatent/.style={draw=black, fill=white, ellipse, inner sep=0.16cm, minimum width=0.42cm, minimum height=0.34cm},
  pmDirected/.style={-{Stealth[length=2mm]}, line width=0.7pt},
  pmBidirected/.style={{Stealth[length=2mm]}-{Stealth[length=2mm]}, line width=0.7pt},
  pmCopath/.style={line width=1.6pt},
  pmLabel/.style={midway, fill=white, inner sep=1pt},
]
  \node[pmLatent] (g_m) at (3.600,9.100) {$g_{m}$};
  \node[pmLatent] (e_m) at (8.200,9.100) {$e_{m}$};
  \node[pmLatent] (g_p) at (16.600,9.100) {$g_{p}$};
  \node[pmLatent] (e_p) at (21.200,9.100) {$e_{p}$};
  \node[pmLatent] (g_o) at (10.100,1.200) {$g_{o}$};
  \node[pmLatent] (e_o) at (5.600,1.200) {$e_{o}$};
  \node[pmLatent] (z_mat_m1) at (0.000,6.200) {$z^{(m)}_{m,1}$};
  \node[pmLatent] (z_mat_m2) at (5.000,6.200) {$z^{(m)}_{m,2}$};
  \node[pmLatent] (z_pat_m1) at (2.200,6.200) {$z^{(p)}_{m,1}$};
  \node[pmLatent] (z_pat_m2) at (7.200,6.200) {$z^{(p)}_{m,2}$};
  \node[pmLatent] (z_mat_p1) at (13.000,6.200) {$z^{(m)}_{p,1}$};
  \node[pmLatent] (z_mat_p2) at (18.000,6.200) {$z^{(m)}_{p,2}$};
  \node[pmLatent] (z_pat_p1) at (15.200,6.200) {$z^{(p)}_{p,1}$};
  \node[pmLatent] (z_pat_p2) at (20.200,6.200) {$z^{(p)}_{p,2}$};
  \node[pmLatent] (z_mat_o1) at (1.100,4.000) {$z^{(m)}_{o,1}$};
  \node[pmLatent] (z_mat_o2) at (6.100,4.000) {$z^{(m)}_{o,2}$};
  \node[pmLatent] (z_pat_o1) at (14.100,4.000) {$z^{(p)}_{o,1}$};
  \node[pmLatent] (z_pat_o2) at (19.100,4.000) {$z^{(p)}_{o,2}$};
  \node[pmObserved] (x_m1) at (1.100,7.700) {$x_{m,1}$};
  \node[pmObserved] (x_m2) at (6.100,7.700) {$x_{m,2}$};
  \node[pmObserved] (y_m) at (3.600,10.500) {$y_{m}$};
  \node[pmObserved] (x_p1) at (14.100,7.700) {$x_{p,1}$};
  \node[pmObserved] (x_p2) at (19.100,7.700) {$x_{p,2}$};
  \node[pmObserved] (y_p) at (16.600,10.500) {$y_{p}$};
  \node[pmObserved] (x_o1) at (7.600,2.600) {$x_{o,1}$};
  \node[pmObserved] (x_o2) at (12.600,2.600) {$x_{o,2}$};
  \node[pmObserved] (y_o) at (10.100,-0.200) {$y_{o}$};
  \draw[pmDirected, draw=black] (z_mat_m1) -- node[pmLabel] {$\frac{1}{2}$} (z_mat_o1);
  \draw[pmDirected, draw=black] (z_pat_m1) -- node[pmLabel] {$\frac{1}{2}$} (z_mat_o1);
  \draw[pmDirected, draw=black] (z_mat_p1) -- node[pmLabel] {$\frac{1}{2}$} (z_pat_o1);
  \draw[pmDirected, draw=black] (z_pat_p1) -- node[pmLabel] {$\frac{1}{2}$} (z_pat_o1);
  \draw[pmDirected, draw=black] (z_mat_m2) -- node[pmLabel] {$\frac{1}{2}$} (z_mat_o2);
  \draw[pmDirected, draw=black] (z_pat_m2) -- node[pmLabel] {$\frac{1}{2}$} (z_mat_o2);
  \draw[pmDirected, draw=black] (z_mat_p2) -- node[pmLabel] {$\frac{1}{2}$} (z_pat_o2);
  \draw[pmDirected, draw=black] (z_pat_p2) -- node[pmLabel] {$\frac{1}{2}$} (z_pat_o2);
  \draw[pmDirected, draw=black] (z_mat_m1) -- (x_m1);
  \draw[pmDirected, draw=black] (z_pat_m1) -- (x_m1);
  \draw[pmDirected, draw=black] (z_mat_m2) -- (x_m2);
  \draw[pmDirected, draw=black] (z_pat_m2) -- (x_m2);
  \draw[pmDirected, draw=black] (x_m1) -- node[pmLabel] {$\beta_{1}$} (g_m);
  \draw[pmDirected, draw=black] (x_m2) -- node[pmLabel] {$\beta_{2}$} (g_m);
  \draw[pmDirected, draw=black] (g_m) -- (y_m);
  \draw[pmDirected, draw=black] (e_m) -- (y_m);
  \draw[pmDirected, draw=black] (z_mat_p1) -- (x_p1);
  \draw[pmDirected, draw=black] (z_pat_p1) -- (x_p1);
  \draw[pmDirected, draw=black] (z_mat_p2) -- (x_p2);
  \draw[pmDirected, draw=black] (z_pat_p2) -- (x_p2);
  \draw[pmDirected, draw=black] (x_p1) -- node[pmLabel] {$\beta_{1}$} (g_p);
  \draw[pmDirected, draw=black] (x_p2) -- node[pmLabel] {$\beta_{2}$} (g_p);
  \draw[pmDirected, draw=black] (g_p) -- (y_p);
  \draw[pmDirected, draw=black] (e_p) -- (y_p);
  \draw[pmDirected, draw=black] (z_mat_o1) -- (x_o1);
  \draw[pmDirected, draw=black] (z_pat_o1) -- (x_o1);
  \draw[pmDirected, draw=black] (z_mat_o2) -- (x_o2);
  \draw[pmDirected, draw=black] (z_pat_o2) -- (x_o2);
  \draw[pmDirected, draw=black] (x_o1) -- node[pmLabel] {$\beta_{1}$} (g_o);
  \draw[pmDirected, draw=black] (x_o2) -- node[pmLabel] {$\beta_{2}$} (g_o);
  \draw[pmDirected, draw=black] (g_o) -- (y_o);
  \draw[pmDirected, draw=black] (e_o) -- (y_o);
  \draw[pmBidirected, draw=black] (z_mat_m1) to[loop above, looseness=6, in=120, out=60] node[pmLabel] {$\frac{1}{2}$} (z_mat_m1);
  \draw[pmBidirected, draw=black] (z_mat_m2) to[loop above, looseness=6, in=120, out=60] node[pmLabel] {$\frac{1}{2}$} (z_mat_m2);
  \draw[pmBidirected, draw=black] (z_pat_m1) to[loop above, looseness=6, in=120, out=60] node[pmLabel] {$\frac{1}{2}$} (z_pat_m1);
  \draw[pmBidirected, draw=black] (z_pat_m2) to[loop above, looseness=6, in=120, out=60] node[pmLabel] {$\frac{1}{2}$} (z_pat_m2);
  \draw[pmBidirected, draw=black] (z_mat_m1) to[bend left=30] (z_pat_m1);
  \draw[pmBidirected, draw=black] (z_mat_m1) to[bend left=30] (z_pat_m2);
  \draw[pmBidirected, draw=black] (z_mat_m2) to[bend left=30] (z_pat_m1);
  \draw[pmBidirected, draw=black] (z_mat_m2) to[bend left=30] (z_pat_m2);
  \draw[pmBidirected, draw=black] (z_mat_p1) to[loop above, looseness=6, in=120, out=60] node[pmLabel] {$\frac{1}{2}$} (z_mat_p1);
  \draw[pmBidirected, draw=black] (z_mat_p2) to[loop above, looseness=6, in=120, out=60] node[pmLabel] {$\frac{1}{2}$} (z_mat_p2);
  \draw[pmBidirected, draw=black] (z_pat_p1) to[loop above, looseness=6, in=120, out=60] node[pmLabel] {$\frac{1}{2}$} (z_pat_p1);
  \draw[pmBidirected, draw=black] (z_pat_p2) to[loop above, looseness=6, in=120, out=60] node[pmLabel] {$\frac{1}{2}$} (z_pat_p2);
  \draw[pmBidirected, draw=black] (z_mat_p1) to[bend left=30] (z_pat_p1);
  \draw[pmBidirected, draw=black] (z_mat_p1) to[bend left=30] (z_pat_p2);
  \draw[pmBidirected, draw=black] (z_mat_p2) to[bend left=30] (z_pat_p1);
  \draw[pmBidirected, draw=black] (z_mat_p2) to[bend left=30] (z_pat_p2);
  \draw[pmBidirected, draw=black] (z_mat_o1) to[loop above, looseness=6, in=120, out=60] node[pmLabel] {$\tfrac14 - \tfrac12 a^{(1)}_{11}$} (z_mat_o1);
  \draw[pmBidirected, draw=black] (z_pat_o1) to[loop above, looseness=6, in=120, out=60] node[pmLabel] {$\tfrac14 - \tfrac12 a^{(1)}_{11}$} (z_pat_o1);
  \draw[pmBidirected, draw=black] (z_mat_o2) to[loop above, looseness=6, in=120, out=60] node[pmLabel] {$\tfrac14 - \tfrac12 a^{(1)}_{22}$} (z_mat_o2);
  \draw[pmBidirected, draw=black] (z_pat_o2) to[loop above, looseness=6, in=120, out=60] node[pmLabel] {$\tfrac14 - \tfrac12 a^{(1)}_{22}$} (z_pat_o2);
  \draw[pmBidirected, draw=black] (e_m) to[loop above, looseness=6, in=120, out=60] node[pmLabel] {$V_{E}$} (e_m);
  \draw[pmBidirected, draw=black] (e_p) to[loop above, looseness=6, in=120, out=60] node[pmLabel] {$V_{E}$} (e_p);
  \draw[pmBidirected, draw=black] (e_o) to[loop above, looseness=6, in=120, out=60] node[pmLabel] {$V_{E}$} (e_o);
  \draw[pmCopath, draw=pmCB03A2E] (y_m) -- node[pmLabel] {$\left[\rho_{y}\right]$} (y_p);
\end{tikzpicture}

        \caption{%
            Two generation-1 individuals mating to produce one generation-2 offspring.
            Unlike Figure~\ref{fig:pair-offspring-2v}, each parent's alleles are correlated.
            No bidirectional arrow joins two alleles on the \emph{same} gamete because $s^{(1)}_{kl} = 0$ by \eqref{eq:s-gen1}.
            The co-path still carries the correlation $\rho_y$, but contributes $\mu^{(1)} = \rho_y/V_Y^{(1)}$, since $V_Y$ has grown.%
        }
        \label{fig:reduced}
    \end{floatcard}
\end{sidewaysfigure}

\paragraph{Covariance across gametes.}
We still think in terms of the four chains described in \eqref{eq:a-gen1}, but now any chain connecting a parent's allele up to its own phenotype is no longer just $\tfrac12\beta_k$: the parent's alleles are correlated now, so a chain leaving $z^{(u)}_{ik}$ can reach the phenotype through \emph{any} other allele the parent carries.
Every such route is already collected in $c^{(t)}_k$, which \eqref{eq:c-def} defined as exactly that covariance.
The four chains then contribute $c^{(t)}_k \cdot \mu^{(t)} \cdot c^{(t)}_l$ in total, so
\begin{equation}\label{eq:a-recursion}
    \boxed{\ a^{(t+1)}_{kl} \;=\; \mut\, c^{(t)}_k\, c^{(t)}_l \ }
\end{equation}
Setting $t=0$ recovers \eqref{eq:a-gen1}, since $c^{(0)}_k = \tfrac12\beta_k$ and $\muo = \rho_y/\VYo$.

\paragraph{Covariance within a gamete.}
When both alleles came from the \emph{same} parent, no chain can leave that parent and therefore the co-path is not crossed:
\begin{equation}\label{eq:s-recursion}
    \boxed{\ s^{(t+1)}_{kl} \;=\; \tfrac14 \Cov[\,x_{ik},\, x_{il}\,] \;=\; \tfrac12\Big(s^{(t)}_{kl} + a^{(t)}_{kl}\Big) \ }
\end{equation}
using \eqref{eq:x-from-as} for the parent's genotype covariance.
In other words, when picking two non-linked alleles from the same gamete in an offspring, they are equally likely to have been from the same vs different gamete in the parent.
Therefore, the covariance between the two offspring's alleles is the average of the across- vs same-gamete covariance in the parent generation.

While $s^{(1)}_{kl} = 0$, the second generation is where the same-gamete covariance becomes non-zero: $s^{(2)}_{kl} = \tfrac12 a^{(1)}_{kl} = \beta_k\beta_l\rho_y/(8\VYo)$.

\paragraph{Covariance between an allele and the phenotype.}
We rewrite the formula for $c$ from \eqref{eq:c-def}, one generation later, as two separate sums:
\begin{equation}\label{eq:c-split}
    c^{(t+1)}_k \;=\; \sum_{l=1}^{M} \beta_l\Big(s^{(t+1)}_{kl} + a^{(t+1)}_{kl}\Big)
    \;=\; \underbrace{\sum_{l} \beta_l\,a^{(t+1)}_{kl}}_{\text{across gametes}}
    \;+\; \underbrace{\sum_{l} \beta_l\,s^{(t+1)}_{kl}}_{\text{same gamete}} .
\end{equation}
The across-gamete part can be written in terms of the previous generation's $c$ by \eqref{eq:a-recursion}.
After pulling out common terms, we are left with a summation that can be simplified according to \eqref{eq:VA-from-as}.
We then plug in $\mut = \rho_y/\VYt$, which leaves $\rho_y\VAt/\VYt$ --- and that is $\rhogt$ by \eqref{eq:rhog-t}:
\begin{equation}\label{eq:c-part-a}
\begin{aligned}
    \sum_{l} \beta_l\,a^{(t+1)}_{kl}
    \;&=\; \sum_{l} \beta_l\,\mut\, c^{(t)}_k c^{(t)}_l
    \;=\; \mut\, c^{(t)}_k \underbrace{\sum_{l} \beta_l c^{(t)}_l}_{\VAt/2}
    \;=\; \frac{\rho_y\,\VAt}{2\,\VYt}\, c^{(t)}_k \\[4pt]
    \;&=\; \tfrac12\,\rhogt\, c^{(t)}_k ,
\end{aligned}
\end{equation}

For the same-gamete part we fill in the recursive expression for $s^{(t+1)}_{kl}$ according to \eqref{eq:s-recursion}, which lets us write it in terms of $c^{(t)}_k$ according to \eqref{eq:c-def}.
However, we must be careful to replace the $k=l$ case with its fixed value of $s^{(t)}_{kk} \equiv \tfrac12$:
\begin{equation}\label{eq:c-part-s}
\begin{aligned}
    \sum_{l} \beta_l\,s^{(t+1)}_{kl}
    &\;=\; \tfrac12 \underbrace{ \sum_{l} \beta_l\big(s^{(t)}_{kl} + a^{(t)}_{kl}\big)}_{c^{(t)}_k}
    \;+\; \beta_k\bigg[\underbrace{\tfrac12}_{s^{(t+1)}_{kk}} \;-\; \underbrace{\tfrac12\big(\tfrac12 + a^{(t)}_{kk}\big)}_{\text{what \eqref{eq:s-recursion} would give for } k=l}\bigg] \\[4pt]
    &\;=\; \tfrac12\,c^{(t)}_k \;+\; \beta_k\Big(\tfrac14 - \tfrac12\,a^{(t)}_{kk}\Big) .
\end{aligned}
\end{equation}
The bracket is the correction for the one term the recursion does not govern.

We now add the two halves.
Note that $\rhogt$ appears in only one of them: it enters through \eqref{eq:c-part-a}, because that is the half whose chains cross the co-path between the parents, and it is absent from \eqref{eq:c-part-s}, whose chains stay inside one parent.
What the two halves share is a factor of $\tfrac12 c^{(t)}_k$, and collecting that factor is what puts $\rhogt$ next to a $1$:
\begin{equation}\label{eq:c-recursion}
\begin{gathered}
    c^{(t+1)}_k \;=\;
    \underbrace{\tfrac12\,\rhogt\, c^{(t)}_k}_{\text{across gametes, \eqref{eq:c-part-a}}}
    \;+\; \underbrace{\tfrac12\, c^{(t)}_k \;+\; \beta_k\Big(\tfrac14 - \tfrac12\,a^{(t)}_{kk}\Big)}_{\text{same gamete, \eqref{eq:c-part-s}}} \\[8pt]
    \boxed{\ c^{(t+1)}_k \;=\; \underbrace{\tfrac12\big(1 + \rhogt\big)\, c^{(t)}_k}_{\text{inherited}}
    \;+\; \beta_k \underbrace{\Big(\tfrac14 - \tfrac12\,a^{(t)}_{kk}\Big)}_{\Var[z^{\perp(m)}_{ok}]} \ }
\end{gathered}
\end{equation}
So the $1$ is the allele the offspring actually received, and the $\rhogt$ beside it is the extra covariance that allele picked up because the parent it came from was chosen to resemble the other parent.
Under random mating $\rhogt = 0$ and only the $1$ remains.
We can notice that the segregation variance \eqref{eq:seg-var-a} naturally emerges from our formula and brings with it a satisfying interpretation.
We have previously thought of offspring alleles as having a directly-inherited component from the parent (equivalent to half the parent's genotype) plus segregation deviation that is uncorrelated with the parent's genotype.
The former contributes half the covariance present in the parent generation, but amplified slightly due to assortative mating predicting the total covariance of correlated variants coming from the other parent.
The latter contributes additional variance directly proportional to that variant's effect size.

We can expand out the terms of equation \eqref{eq:c-recursion} such that they only rely on the starting parameters --- the effect sizes $\beta_k$, the phenotypic correlation between mates $\rho_y$, and the environmental variance $V_E$ --- and previous states of $c$.
Substituting $\rhogt = \rho_y\VAt/\VYt$, then $\VAt = 2\sum_j \beta_j c^{(t)}_j$ and $\VYt = \VAt + V_E$, and finally $a^{(t)}_{kk} = \mu^{(t-1)}\big(c^{(t-1)}_k\big)^2$ leaves

\begin{equation}\label{eq:c-closed}
    c^{(t+1)}_k \;=\; \tfrac14\beta_k
    \;+\; \tfrac12\left(1 + \frac{2\rho_y\sum_j \beta_j c^{(t)}_j}{2\sum_j \beta_j c^{(t)}_j + V_E}\right) c^{(t)}_k
    \;-\; \frac{\rho_y\,\beta_k}{2\left(2\sum_j \beta_j c^{(t-1)}_j + V_E\right)} \big(c^{(t-1)}_k\big)^2 ,
\end{equation}
with $c^{(t)}_k = \tfrac12\beta_k$ for every $t \leq 0$, the base population having been reached without assortment.

This reveals that the covariances in this population follow an \emph{autonomous second-order difference equation}: autonomous because the update rule is the same at every generation, and second-order because it reaches back two of them.
The $M$ variants form one coupled system rather than $M$ separate equations, since each variant's update reads all the others through $\VAt$.
It is also \emph{nonlinear} --- $c^{(t-1)}_k$ appears squared, and $c^{(t)}$ appears inside a denominator --- so the usual closed-form solutions for linear recursions do not apply, and we look instead for a \emph{fixed point}: a $c$ that the update leaves unchanged.
The leading factor $\tfrac12\big(1 + \rhogt\big)$ is less than one whenever $\rhogt < 1$, which makes the update a contraction near such a point, so the approach to it is geometric rather than oscillatory or divergent: the remaining distance falls by a roughly constant fraction each generation.

\begin{table}[H]
    \begin{floatcard}
    \renewcommand{\arraystretch}{2.1}
    \setlength{\tabcolsep}{2pt}
    \footnotesize
    \begin{tabular}{ll|ccc}
            & & $t \leq 0$ \ \footnotesize base pop. & $t = 1$ & general $(t)$ \\
        \hline
        allele--phenotype & $c^{(t)}_k$
            & $\tfrac12\beta_k$
            & $\tfrac12\beta_k\big(1 + \tfrac12\rhogo\big)$
            & $\begin{aligned}[c]
                   \tfrac12&\big(1 + \rho_g^{(t-1)}\big) c^{(t-1)}_k \\[-2pt]
                   &+\; \beta_k\Big(\tfrac14 - \tfrac12 a^{(t-1)}_{kk}\Big)
               \end{aligned}$ \\
        across gametes & $a^{(t)}_{kl}$
            & $0$
            & $\dfrac{\beta_k\beta_l}{4\,\VAo}\rhogo$
            & $\mu^{(t-1)}\, c^{(t-1)}_k c^{(t-1)}_l$ \\
        same gamete & $s^{(t)}_{kl}$, $k \neq l$
            & $0$
            & $0$
            & $\tfrac12\big(s^{(t-1)}_{kl} + a^{(t-1)}_{kl}\big)$ \\
        \hline
        genotype variance & $\Var[x_{ik}]$
            & $1$
            & $1 + \dfrac{\beta_k^2}{2\,\VAo}\rhogo$
            & $1 + 2\,a^{(t)}_{kk}$ \\
        segregation variance & $\Var\big[z^{\perp(m)}_{ok}\big]$
            & $\tfrac14$
            & $\tfrac14 - \dfrac{\beta_k^2}{8\,\VAo}\rhogo$
            & $\tfrac14 - \tfrac12 a^{(t)}_{kk}$ \\
        additive variance & $\VAt$
            & $\VAo$
            & $\VAo\big(1 + \tfrac12\rhogo\big)$
            & $2\sum_k \beta_k c^{(t)}_k$ \\
        phenotypic variance & $\VYt$
            & $\VAo + V_E$
            & $V_A^{(1)} + V_E$
            & $\VAt + V_E$ \\
        heritability & $\hsqt$
            & $\VAo/\VYo$
            & $V_A^{(1)}/V_Y^{(1)}$
            & $\VAt/\VYt$ \\
        genetic correlation & $\rhogt$
            & $\rho_y\,\hsqo$
            & $\rho_y\,h^2_{(1)}$
            & $\rho_y\,\hsqt$ \\
        co-path coefficient & $\mut$
            & $\rho_y/\VYo$
            & $\rho_y/V_Y^{(1)}$
            & $\rho_y/\VYt$ \\
    \end{tabular}
    \caption{%
        Every quantity derived in this section, from the base population forward, for a trait with $M$ variants.
        Every column gives that generation's own value, the left-hand side always being the quantity named in the second column.
        The three recursions therefore reach back to $t-1$ in the last column, while the rest are evaluated at $t$ directly.
        The first block are the three allele covariances, of which only $c$ is a recursion the others do not feed: $a$ and $s$ are read off it, and $s^{(t)}_{kk} \equiv \tfrac12$.
        The second block is read off the same generation's $c$, and involves no recursion at all.
        The base-population column is written $t \leq 0$ because nothing has happened yet: its ancestors mated at random, so $c$, $a$ and $s$ take these values at every $t \leq 0$, which is what lets the second-order \eqref{eq:c-closed} start.
        Assortment begins with the generation-0 pairing, so $\rhogt$ and $\mut$ in that column are the ones governing the step to $t=1$.%
    }
    \label{tab:as-recursion}
    \end{floatcard}
\end{table}


\subsubsection{Equilibrium}
Our next goal is to solve for the values that lead to an equilibrium state in the recursive covariance formulas we expressed above.
An equilibrium is a state the update in \eqref{eq:c-recursion} leaves unchanged, so we look for a $c$ satisfying $c^{(t+1)}_k = c^{(t)}_k = c^{(t-1)}_k$, which we write $c^{(\mathrm{eq})}_k$.

\begin{assumption}{AM-equilibrium}{Intergenerational equilibrium}
Assortment has been operating at the same strength for long enough that the covariances have stopped changing from one generation to the next.
\bundles{equilibrium assortative mating; reached after roughly six to ten generations \citep{Sunde2024NatComm}.}
\end{assumption}

Under equilibrium, we fill in the same value, $c_k^{(eq)}$ for all states of $c_k$ in \eqref{eq:c-recursion} after replacing $a^{(eq)}_{kk} = \mu^{(eq)}\big(c^{(eq)}_k\big)^2$.
\begin{equation}\label{eq:c-eq-quadratic}
\begin{aligned}
    c^{(\mathrm{eq})}_k \;&=\; \tfrac12\big(1+\rhogeq\big)c^{(\mathrm{eq})}_k
    \;+\; \beta_k\Big(\tfrac14 - \tfrac12\,\mueq \big(c^{(\mathrm{eq})}_k\big)^2\Big) \\[4pt]
    \tfrac12\big(1-\rhogeq\big)\,c^{(\mathrm{eq})}_k \;&=\; \tfrac14\beta_k
    \;-\; \tfrac12\,\mueq\,\beta_k \big(c^{(\mathrm{eq})}_k\big)^2 \\[4pt]
    2\big(1-\rhogeq\big)\,c^{(\mathrm{eq})}_k \;&=\; \beta_k
    \;-\; 2\,\mueq\,\beta_k \big(c^{(\mathrm{eq})}_k\big)^2 \\[4pt]
    2\mueq\beta_k \big(c^{(\mathrm{eq})}_k\big)^2
    + 2\big(1-\rhogeq\big) c^{(\mathrm{eq})}_k - \beta_k \;&=\; 0
\end{aligned}
\end{equation}
\begin{equation}\label{eq:c-eq}
\begin{gathered}
    c^{(\mathrm{eq})}_k \;=\; \frac{-\big(1-\rhogeq\big) \pm \sqrt{\big(1-\rhogeq\big)^2 + 2\,\mueq\,\beta_k^2}}{2\,\mueq\,\beta_k} \\[8pt]
    \boxed{\, c^{(\mathrm{eq})}_k \,=\, \frac{\beta_k}{\big(1-\rhogeq\big) + \sqrt{\big(1-\rhogeq\big)^2 + 2\mueq\beta_k^2}} \,}
\end{gathered}
\end{equation}
Only the $+$ root is admissible: the $-$ root gives a $c^{(\mathrm{eq})}_k$ of the opposite sign to $\beta_k$, and does not reduce to the base-population value when assortment is switched off.

However, we still have the problem that $\mueq$ and $\rhogeq$ depend on $\VAeq$, which itself depends on the sum of $c^{(eq)}_k$ for all $k$, which are quantities we are solving for.
Summing \eqref{eq:c-eq} over variants does not resolve this, because each variant carries its own square root containing its own $2\mueq\beta_k^2$, so the $M$ terms cannot be cleanly gathered into one.
To make obtaining a closed form solution manageable, we need to apply an approximation that is justified by the following assumptions:

\begin{assumption}{Large-Me}{The additive variance is spread across many variants}
Taking inspiration from \citet{OConnor2019AJHG}, define the \emph{effective number of variants} as
\begin{equation}\label{eq:Me-def}
    \boxed{\ M_e \;:=\; \frac{\big(\VAo\big)^2}{\sum_k \beta_k^4} \;=\; \frac{\big(\sum_k \beta_k^2\big)^2}{\sum_k \beta_k^4} \ }
\end{equation}
This is closely related to the same quantity defined in \citet{OConnor2019AJHG}, albeit ours is smaller by a factor of three because we are aiming for the property that $M_e = M$ when all variants have the same effect size, rather than when all variants have the same \emph{Gaussian prior}.

Specifically, we assume that the effective number of variants is large:
$$M_e \;\gg\; 1 .$$

Equivalently, no single variant explains an appreciable share of the additive variance.
Since $\sum_j \beta_j^4 \geq \beta_k^4$ for any single $k$, rearranging \eqref{eq:Me-def} gives
\begin{equation}\label{eq:Me-bound}
    \frac{\beta_k^2}{\VAo} \;\leq\; \frac{1}{\sqrt{M_e}} \qquad \text{for every } k ,
\end{equation}
so $M_e \gg 1$ forces every variance share below $M_e^{-1/2} \ll 1$.
\bundles{the polygenic, or infinitesimal, model with large $M$; $M_e$ is estimable from summary statistics \citep{OConnor2019AJHG}.}
\end{assumption}

\begin{assumption}{Moderate-AM}{Assortment is not maximal}
The genetic correlation between mates stays bounded away from its ceiling, $\rhogeq < 1$, so that $1 - \rhogeq$ is not close to zero.
\bundles{any trait that is not both highly heritable and highly assorted.}
\end{assumption}

\begin{approximation}{Uniform-inflation}{Assortment inflates every variant's covariance with the phenotype alike}
\begin{equation}\label{eq:c-eq-approx}
\begin{gathered}
    c^{(\mathrm{eq})}_k \;=\; \frac{\beta_k}{\big(1-\rhogeq\big)
    + \sqrt{\big(1-\rhogeq\big)^2 + \underbrace{2\,\mueq\,\beta_k^2}_{\text{negligible}}}} \\[8pt]
    \boxed{\ c^{(\mathrm{eq})}_k \;\approx\; \frac{\beta_k}{2\big(1-\rhogeq\big)}\ }
\end{gathered}
\end{equation}
Every variant is therefore inflated above its base-population value $\tfrac12\beta_k$ by the same factor $1/(1-\rhogeq)$, whatever its effect size.
\justification{The neglected term, relative to the one it sits beside, is
$$\frac{2\,\mueq\,\beta_k^2}{\big(1-\rhogeq\big)^2}
\;=\; \underbrace{\frac{2\,\rhogeq}{1-\rhogeq}}_{\text{\assumptionid{Moderate-AM}}} \cdot \underbrace{\frac{\beta_k^2}{\VAo}}_{\text{\assumptionid{Large-Me}}} ,$$
using $\mueq = \rhogeq/\VAeq$ together with the $\VAeq = \VAo/(1-\rhogeq)$ that this approximation itself yields below in \eqref{eq:VA-eq}.
The first factor is bounded by \assumptionid{Moderate-AM} and the second is small by \assumptionid{Large-Me} via \eqref{eq:Me-bound}, so neither alone suffices.
All that's required is that the product of the two is close to 0.}
\end{approximation}

We'll make use of this approximation to obtain a clean formula for $\VAeq$, but first, we derive the exact formula.
Expanding $\Var[g_i]$ over all pairs of variants, every off-diagonal covariance is $4a^{(\mathrm{eq})}_{kl}$ --- that is \eqref{eq:x-from-as} with $s^{(\mathrm{eq})}_{kl} = a^{(\mathrm{eq})}_{kl}$, the fixed point of \eqref{eq:s-recursion} --- while the diagonal is instead $1 + 2a^{(\mathrm{eq})}_{kk}$ by \eqref{eq:varx-a}.
We apply the off-diagonal rule to every pair and then correct the one term it does not govern, using
$a^{(\mathrm{eq})}_{kl} = \mueq c^{(\mathrm{eq})}_k c^{(\mathrm{eq})}_l$ from \eqref{eq:a-recursion},
$\mueq\VAeq = \rhogeq$ from \eqref{eq:rhog-t},
and $\sum_k \beta_k c^{(\mathrm{eq})}_k = \tfrac12\VAeq$ from \eqref{eq:VA-from-as}:
\begin{equation}\label{eq:VA-eq-exact}
\begin{gathered}
\begin{aligned}
    \VAeq \;&=\; \sum_k\sum_l \beta_k\beta_l \Cov[\,x_{ik},\, x_{il}\,] \\[4pt]
    \;&=\; \underbrace{\sum_k\sum_l \beta_k\beta_l\,4a^{(\mathrm{eq})}_{kl}}_{\text{one rule for every pair}}
    \;+\; \underbrace{\sum_k \beta_k^2\Big[\big(1 + 2a^{(\mathrm{eq})}_{kk}\big) - 4a^{(\mathrm{eq})}_{kk}\Big]}_{\text{diagonal correction}} \\[4pt]
    \;&=\; 4\,\mueq\Big(\sum_k \beta_k c^{(\mathrm{eq})}_k\Big)^{\!2}
    \;+\; \sum_k \beta_k^2\big(1 - 2a^{(\mathrm{eq})}_{kk}\big) \\[4pt]
    \;&=\; \underbrace{\mueq\big(\VAeq\big)^2}_{\rhogeq\,\VAeq}
    \;+\; \VAo \;-\; 2\sum_k \beta_k^2 a^{(\mathrm{eq})}_{kk} \\[4pt]
    \VAeq\big(1-\rhogeq\big) \;&=\; \VAo \;-\; 2\sum_k \beta_k^2 a^{(\mathrm{eq})}_{kk}
\end{aligned} \\[8pt]
    \boxed{\ \VAeq \;=\; \frac{\VAo - 2\sum_k \beta_k^2\, a^{(\mathrm{eq})}_{kk}}{1-\rhogeq}\ }
\end{gathered}
\end{equation}

We'll see in a later section that half of this numerator, $\tfrac12\big(\VAo - 2\sum_k \beta_k^2 a^{(\mathrm{eq})}_{kk}\big)$, will come up again as the variance of the deviation of offspring genetic values from their parents' average.
We use \eqref{eq:a-recursion} to write $a$ in terms of $c$ and then plug in the approximated formula for $c$ from \eqref{eq:c-eq-approx}:
\begin{equation}\label{eq:VA-eq-drag}
\begin{aligned}
    2\sum_k \beta_k^2\, a^{(\mathrm{eq})}_{kk}
    \;&=\; 2\,\mueq \sum_k \beta_k^2 \big(c^{(\mathrm{eq})}_k\big)^2
    \;\approx\; 2\,\mueq \sum_k \beta_k^2\,\frac{\beta_k^2}{4\big(1-\rhogeq\big)^2} \\[4pt]
    \;&=\; \frac{\mueq}{2\big(1-\rhogeq\big)^2}\sum_k \beta_k^4 \\[4pt]
    \;&=\; \frac{\rhogeq}{2\,\VAo\big(1-\rhogeq\big)}\sum_k \beta_k^4 ,
\end{aligned}
\end{equation}
the last step using $\mueq = \rhogeq\big(1-\rhogeq\big)/\VAo$, which is the exact $\mueq = \rhogeq/\VAeq$ with $\VAeq$ replaced by the very approximation we are in the middle of deriving.
This justified because this value of $\VAeq$ is used only to \emph{measure the size} of a term we are about to discard, and the error in this discarded term is guaranteed to be smaller than the error in the term we are approximating.

Substituting that back into \eqref{eq:VA-eq-exact} and factoring out $\VAo/\big(1-\rhogeq\big)$,
\begin{equation}\label{eq:VA-eq-correction}
    \VAeq \;\approx\; \frac{\VAo}{1-\rhogeq}\left[1 \;-\; \frac{\rhogeq}{2\big(1-\rhogeq\big)} \cdot \frac{\sum_k \beta_k^4}{\big(\VAo\big)^2}\right]
\end{equation}
The last factor of \eqref{eq:VA-eq-correction} is exactly $1/M_e$ by \eqref{eq:Me-def}, so the whole correction is $\rhogeq/\big[2(1-\rhogeq)M_e\big]$ --- an error controlled by the effective number of variants rather than by any individual effect size.

\begin{approximation}{VA-inflation}{The additive variance inflates by a single factor}
The bracketed correction in \eqref{eq:VA-eq-correction} is negligible, leaving the additive variance at equilibrium as the base-population value scaled by one factor that depends on nothing but the genetic correlation between mates:
\begin{equation}\label{eq:VA-eq}
    \boxed{\ \VAeq \;=\; \frac{\VAo}{1-\rhogeq}\ }
\end{equation}
\justification{Written with \eqref{eq:Me-def}, the neglected term is a relative error in $\VAeq$ of
$$\frac{\rhogeq}{2\big(1-\rhogeq\big)} \cdot \frac{1}{M_e} ,$$
small by \assumptionid{Large-Me} and bounded by \assumptionid{Moderate-AM}.
The error is signed: dropping the term \emph{overstates} $\VAeq$, since a variant's own departure from Hardy--Weinberg drags its contribution down.}
\end{approximation}

We can also write a formula for the equilibrium heritability:
\begin{equation}\label{eq:hsq-eq}
\begin{gathered}
    \hsqeq \;=\; \frac{\VAeq}{\VYeq} \;=\; \frac{\VAeq}{\VAeq + V_E}
    \;=\; \frac{\VAo/\big(1-\rhogeq\big)}{\VAo/\big(1-\rhogeq\big) + V_E} \\[8pt]
    \boxed{\ \hsqeq \;=\; \frac{\VAo}{\VAo + \big(1-\rhogeq\big)V_E}\ } \\[8pt]
    \hsqeq \;=\; \frac{\hsqo}{\hsqo + \big(1-\rhogeq\big)\big(1-\hsqo\big)} \\[8pt]
    \boxed{\ \hsqeq \;=\; \frac{\hsqo}{1 - \rhogeq\big(1-\hsqo\big)}\ }
\end{gathered}
\end{equation}

These formulas for $\VAeq$ and $\hsqeq$ are easily interpretable, but $\rhogeq$ still relies on either $\VAeq$ or $\hsqeq$ through \eqref{eq:rhog-t}.
We can define $\rhogeq$ purely in terms of the starting parameters to resolve the circularity.

\begin{equation}\label{eq:rhog-quadratic}
\begin{aligned}
    \rhogeq \;&=\; \rho_y\,\hsqeq \;=\; \frac{\rho_y\,\hsqo}{1 - \rhogeq\big(1-\hsqo\big)} \\[4pt]
    \rhogeq \;-\; \big(\rhogeq\big)^2\big(1-\hsqo\big) \;&=\; \rho_y\,\hsqo \\[4pt]
    \big(1 - \hsqo\big)\big(\rhogeq\big)^2 \;-\; \rhogeq \;+\; \rho_y\,\hsqo \;&=\; 0
\end{aligned}
\end{equation}
\begin{equation}\label{eq:rhog-eq}
    \boxed{\ \rhogeq \;=\; \frac{1 - \sqrt{1 - 4\,\rho_y\,\hsqo\big(1-\hsqo\big)}}{2\big(1-\hsqo\big)} \ }
\end{equation}
The $-$ root is the admissible one and the discriminant is never negative.

\subsection{Covariance between relatives}
Now that we have derived the statistical properties of the population at equilibrium under assortative mating, we will determine the expected covariance between pairs of individuals given some set of information, similar to Section~\ref{sec:random-mating} of this write-up.

\subsubsection{Inheritance of genetic values}
Rather than modelling the inheritance of allele states from parent to offspring, we will instead model the transmission of genetic values themselves.
From here on we work only at equilibrium, and drop the $(\mathrm{eq})$ superscript for legibility: $V_A$, $V_Y$, $h^2$, $\rho_g$ and $\mu$ now denote the equilibrium values solved for above.

We can extend the logic of \eqref{eq:transmit}--\eqref{eq:seg-def} to decompose the genetic value of an offspring as being the average of their parents' genetic values plus an independent residual term.
Summing the transmitted alleles over variants, weighted by their effects, and substituting \eqref{eq:seg-def} for each one:

\begin{equation}\label{eq:g-transmit}
\begin{gathered}
\begin{aligned}
    g_o \;&=\; \sum_k \beta_k\, x_{ok} \;=\; \sum_k \beta_k\big(z^{(m)}_{ok} + z^{(p)}_{ok}\big) \\[4pt]
    \;&=\; \sum_k \beta_k\Big(\tfrac12 x_{mk} + z^{\perp(m)}_{ok} \;+\; \tfrac12 x_{pk} + z^{\perp(p)}_{ok}\Big) \\[4pt]
    \;&=\; \tfrac12 g_m \;+\; \tfrac12 g_p
    \;+\; \underbrace{\sum_k \beta_k\big(z^{\perp(m)}_{ok} + z^{\perp(p)}_{ok}\big)}_{\textstyle g^{\perp}_o}
\end{aligned} \\[8pt]
    \boxed{\ g_o \;=\; \tfrac12 g_m \;+\; \tfrac12 g_p \;+\; g^{\perp}_o \ }
\end{gathered}
\end{equation}

Since each $z^{\perp(u)}_{ok}$ is uncorrelated with \emph{every} allele carried by the parent it came from, their weighted sum $g^{\perp}_o$ is also going to be uncorrelated with $g_m$ and $g_p$ and everything upstream of it.
As a consequence, we can easily compute the variance of this genetic value deviation, making use of \eqref{eq:seg-var-a}:
\begin{equation}\label{eq:g-seg-var}
\begin{gathered}
\begin{aligned}
    \Var[g^{\perp}_o] \;&=\; \sum_k\sum_l \beta_k\beta_l \Cov\!\big[z^{\perp(m)}_{ok} + z^{\perp(p)}_{ok},\; z^{\perp(m)}_{ol} + z^{\perp(p)}_{ol}\big] \\[4pt]
    \;&=\; \sum_k \beta_k^2\Big(\Var\!\big[z^{\perp(m)}_{ok}\big] + \Var\!\big[z^{\perp(p)}_{ok}\big]\Big)
    \;=\; \sum_k \beta_k^2\big(\tfrac12 - a_{kk}\big)
\end{aligned} \\[8pt]
    \boxed{\ \Var[g^{\perp}_o] \;=\; \tfrac12 V_A\big(1 - \rho_g\big) \approx \tfrac12 \VAo \ }
\end{gathered}
\end{equation}
Note that under \approximationid{VA-inflation}, which is \eqref{eq:VA-eq} rearranged to $V_A\big(1-\rho_g\big) = \VAo$, this is equivalent to half of the additive variance under random mating.
This means that the residual genetic variance of a child around the parental mean is determined by the additive variance under random mating and therefore is not impacted by assortative mating.

Because $g^{\perp}_o$ is uncorrelated with both parents and across siblings, \eqref{eq:g-transmit} is a linear model of exactly the kind Section~\ref{sec:path-analysis} draws: two directed edges of $\tfrac12$ into the offspring's genetic value, and a disturbance carrying \eqref{eq:g-seg-var}.
We can therefore draw a simplified path diagram that omits individual alleles and genotypes.
We'll also omit the explicit environmental term and instead code it as exogeneous variance on the phenotype.

\begin{figure}[H]
    \begin{floatcard}
        \definecolor{pmCB03A2E}{HTML}{B03A2E}
\begin{tikzpicture}[
  every node/.style={font=\small},
  pmObserved/.style={draw=black, fill=white, rectangle, inner sep=0.09cm, minimum width=0.42cm, minimum height=0.34cm},
  pmLatent/.style={draw=black, fill=white, ellipse, inner sep=0.16cm, minimum width=0.42cm, minimum height=0.34cm},
  pmDirected/.style={-{Stealth[length=2mm]}, line width=0.7pt},
  pmBidirected/.style={{Stealth[length=2mm]}-{Stealth[length=2mm]}, line width=0.7pt},
  pmCopath/.style={line width=1.6pt},
  pmLabel/.style={midway, fill=white, inner sep=1pt},
]
  \node[pmLatent] (g_m) at (2.000,3.400) {$g_m$};
  \node[pmLatent] (g_p) at (8.000,3.400) {$g_p$};
  \node[pmLatent] (g_o1) at (3.300,1.600) {$g_{o_1}$};
  \node[pmLatent] (g_o2) at (6.700,1.600) {$g_{o_2}$};
  \node[pmObserved] (y_m) at (0.600,5.200) {$y_m$};
  \node[pmObserved] (y_p) at (9.400,5.200) {$y_p$};
  \node[pmObserved] (y_o1) at (2.100,0.000) {$y_{o_1}$};
  \node[pmObserved] (y_o2) at (7.900,0.000) {$y_{o_2}$};
  \draw[pmDirected, draw=black] (g_m) -- (g_o1);
  \draw[pmDirected, draw=black] (g_p) -- (g_o1);
  \draw[pmDirected, draw=black] (g_m) -- (g_o2);
  \draw[pmDirected, draw=black] (g_p) -- (g_o2);
  \draw[pmDirected, draw=black] (g_m) -- (y_m);
  \draw[pmDirected, draw=black] (g_p) -- (y_p);
  \draw[pmDirected, draw=black] (g_o1) -- (y_o1);
  \draw[pmDirected, draw=black] (g_o2) -- (y_o2);
  \draw[pmBidirected, draw=black] (g_m) to[loop above, looseness=6, in=120, out=60] node[pmLabel] {$V_{A}$} (g_m);
  \draw[pmBidirected, draw=black] (g_p) to[loop above, looseness=6, in=120, out=60] node[pmLabel] {$V_{A}$} (g_p);
  \draw[pmBidirected, draw=black] (g_o1) to[loop above, looseness=6, in=120, out=60] node[pmLabel] {$\tfrac12 V_A\big(1-\rho_g\big)$} (g_o1);
  \draw[pmBidirected, draw=black] (g_o2) to[loop above, looseness=6, in=120, out=60] node[pmLabel] {$\tfrac12 V_A\big(1-\rho_g\big)$} (g_o2);
  \draw[pmBidirected, draw=black] (y_m) to[loop above, looseness=6, in=120, out=60] node[pmLabel] {$V_{E}$} (y_m);
  \draw[pmBidirected, draw=black] (y_p) to[loop above, looseness=6, in=120, out=60] node[pmLabel] {$V_{E}$} (y_p);
  \draw[pmBidirected, draw=black] (y_o1) to[loop above, looseness=6, in=120, out=60] node[pmLabel] {$V_{E}$} (y_o1);
  \draw[pmBidirected, draw=black] (y_o2) to[loop above, looseness=6, in=120, out=60] node[pmLabel] {$V_{E}$} (y_o2);
  \draw[pmCopath, draw=pmCB03A2E] (y_m) -- node[pmLabel] {$\left[\rho_{y}\right]$} (y_p);
\end{tikzpicture}

        \caption{%
            A mated pair and two of their offspring at equilibrium, with inheritance happening at the level of genetic values, rather than alleles.
            The environment appears only as the $V_E$ on each phenotype's own self-loop, so $\Var[y] = V_A + V_E$ as before.%
        }
        \label{fig:g-only}
    \end{floatcard}
\end{figure}

Similar to Table~\ref{tab:relatives-gen1}, we can summarize the covariance between the genetic and phenotypic values across mate pairs and across parent--offspring pairs, now at equilibrium rather than in the first generation.
These expressions are exact if we hold \assumptionid{No-inbreeding}, \assumptionid{Discrete-generations}, and \assumptionid{AM-equilibrium}.

\begin{table}[H]
    \begin{floatcard}
    \renewcommand{\arraystretch}{1.9}
    \setlength{\tabcolsep}{6pt}
    \footnotesize
    \begin{tabular}{l|cc}
            & mate pair $(m,p)$
            & parent--offspring $(m,o_1)$ \\
        \hline
        $\Cov[\,g_m,\; g_\ast\,]$
            & $\rho_g V_A$
            & $\tfrac12 V_A\big(1+\rho_g\big)$ \\
        $\Cov[\,g_m,\; y_\ast\,]$
            & $\rho_y V_A$
            & $\tfrac12 V_A\big(1+\rho_g\big)$ \\
        $\Cov[\,y_m,\; g_\ast\,]$
            & $\rho_y V_A$
            & $\tfrac12 V_A\big(1+\rho_y\big)$ \\
        $\Cov[\,y_m,\; y_\ast\,]$
            & $\rho_y V_Y$
            & $\tfrac12 V_A\big(1+\rho_y\big)$ \\
    \end{tabular}
    \caption{%
        Covariances read off Figure~\ref{fig:g-only} at equilibrium, where $\ast$ stands for the second member of the pair named by the column.%
    }
    \label{tab:g-only-cov}
    \end{floatcard}
\end{table}

Note the asymmetry that arises:
The covariance between a parent's genetic value and their offspring's phenotype is lower than the covariance between a parent's phenotypic value and their offspring's genetic value.
This is because the parent's environment (which is part of their phenotype) predicts the genetic value their offspring inherits from the other parent but not the offspring's environmental value, which remains random under our current model.

% --- HIDDEN (not rendered): the mean-parent abstraction. Nothing currently cites it, but it is
% the natural language for a couple acting as one unit, so it is kept in source to revisit;
% unhide by removing the \iffalse ... \fi wrapper. ---
\iffalse

We can also conceive of an abstract ``mean-parent'' that averages the genetic, environmental, and phenotypic values of a pair of mates:
\begin{equation}\label{eq:mean-parent}
    g_{mp} \;=\; \tfrac12\big(g_m + g_p\big), \qquad
    e_{mp} \;=\; \tfrac12\big(e_m + e_p\big), \qquad
    y_{mp} \;=\; \tfrac12\big(y_m + y_p\big) \;=\; g_{mp} + e_{mp}
\end{equation}
We can obtain the variance of the mean-parent from Table~\ref{tab:g-only-cov}:
\begin{equation}\label{eq:mean-parent-var}
\begin{gathered}
\begin{aligned}
    \Var[g_{mp}] \;&=\; \tfrac14\Big(\Var[g_m] + \Var[g_p] + 2\Cov[\,g_m,\, g_p\,]\Big)
    \;=\; \tfrac14\big(2V_A + 2\rho_g V_A\big) \\[4pt]
    \Var[y_{mp}] \;&=\; \tfrac14\Big(\Var[y_m] + \Var[y_p] + 2\Cov[\,y_m,\, y_p\,]\Big)
    \;=\; \tfrac14\big(2V_Y + 2\rho_y V_Y\big)
\end{aligned} \\[8pt]
    \boxed{\ \Var[g_{mp}] \;=\; \tfrac12 V_A\big(1+\rho_g\big)\ } \qquad\qquad
    \boxed{\ \Var[y_{mp}] \;=\; \tfrac12 V_Y\big(1+\rho_y\big)\ }
\end{gathered}
\end{equation}

We can then encode the genetic value of the offspring as being a linear function of this mean-parent's genetic value, which collapses \eqref{eq:g-transmit} to a single edge:
\begin{equation}\label{eq:g-transmit-mp}
    \boxed{\ g_o \;=\; g_{mp} \;+\; g^{\perp}_o \ }
\end{equation}
The residual $g^{\perp}_o$ still carries \eqref{eq:g-seg-var} and is still uncorrelated with both parents, hence with $g_{mp}$, so the two terms partition the offspring's genetic variance exactly:
\begin{equation}\label{eq:g-partition}
    \Var[g_o] \;=\; \Var[g_{mp}] + \Var[g^{\perp}_o]
    \;=\; \tfrac12 V_A\big(1+\rho_g\big) + \tfrac12 V_A\big(1-\rho_g\big) \;=\; V_A
\end{equation}

We can then similarly write a table for the covariance in genetic and phenotypic values of offspring to their own mean-parent.

\begin{table}[H]
    \begin{floatcard}
    \renewcommand{\arraystretch}{1.9}
    \setlength{\tabcolsep}{10pt}
    \footnotesize
    \begin{tabular}{l|ccc}
            & $\Var[\,\cdot\,]$
            & $\Cov[\,\cdot,\; g_o\,]$
            & $\Cov[\,\cdot,\; y_o\,]$ \\
        \hline
        $\cdot = g_{mp}$
            & $\tfrac12 V_A\big(1+\rho_g\big)$
            & $\tfrac12 V_A\big(1+\rho_g\big)$
            & $\tfrac12 V_A\big(1+\rho_g\big)$ \\
        $\cdot = y_{mp}$
            & $\tfrac12 V_Y\big(1+\rho_y\big)$
            & $\tfrac12 V_A\big(1+\rho_y\big)$
            & $\tfrac12 V_A\big(1+\rho_y\big)$ \\
    \end{tabular}
    \caption{%
        Variances of a mean-parent and its covariances with one of its offspring, at equilibrium.
        The two right-hand columns reproduce the parent--offspring column of Table~\ref{tab:g-only-cov}, since averaging two exchangeable parents returns the same covariance of each parent with the child they share.%
    }
    \label{tab:mean-parent-cov}
    \end{floatcard}
\end{table}

\fi
% --- end hidden segment ---

\subsubsection{Framework for deriving covariance formulas}
Our next goal is to derive formulas that express the covariance between phenotypes and/or genetic values of any two individuals $X$ and $Y$ in a pedigree.
We adopt graph theory notation and terminology to analyze their relationship.
We read the pedigree as a \emph{graph}: a set of \emph{nodes}, one for each individual, joined by two types of \emph{edges}:

\begin{itemize}
    \item An \emph{undirected} edge, written $X \,\text{---}\, Y$, for each \textbf{mating}.
    \item A \emph{directed} edge, written $X \rightarrow Y$, for each \textbf{meiosis}, pointing from parent to offspring.
\end{itemize}


    A \emph{mating graph} consists of only the mating edges, while a \emph{meiosis graph} consists of only the meiosis edges.
    We will be defining \emph{mating chains} in the former, and \emph{lineal chains} in the latter.
    Then, a path $\mathcal{P}$ between two probands in a pedigree becomes an alternating sequence of mating and lineal chains, and properties of that sequence and its chains will be enough to write the covariance between the probands in closed form.


Up to this point, discussion of the genetic consequences of assortative mating requires the definition of mates to be strictly parents that reproduced.
However, for the rest of this section, the definition of mates is simply any primary assortment based on the phenotype that occurs in a population with the same strength as that of assortartive mating $\rho_y$ between parents.
Therefore, this can include homosexual pairs and heterosexual pairs that do not reproduce.
In graph terms: the mating graph gets an edge for every such pair, whether or not a meiosis edge leaves it.

\subsubsection{Mating chains}
The \emph{mating graph} consists of only the undirected mating edges between nodes.
A \emph{path} in it is a sequence of distinct nodes, each joined by an edge to the next.
The following assumption makes working with mating graphs feasible for our purposes:


\begin{assumption}{No-mating-loops}{The mating graph is a forest}
    A \emph{cycle} in the mating graph is a path whose last node is joined by an edge back to its first.
    A graph with no cycles is a \emph{forest}, and each of its connected components is a \emph{tree}.

    We assume the mating graph is a forest.
    As a consequence, any two nodes are joined at most by one path, and therefore must be on the same tree.
\end{assumption}

For $X \neq Y$ in the same tree, the \textbf{mating chain} $\Mu(X,Y)$ is the unique path joining them, read as an ordered list running from $X$ to $Y$.
The \textbf{mating distance} $\Nmu(X,Y)$ is the number of edges on that path, i.e. the number of matings crossed in getting from $X$ to $Y$, and is one fewer than the size of it:
\begin{equation}\label{eq:mating-dist}
    \boxed{\ \Nmu(X,Y) \;=\; \big|\Mu(X,Y)\big| - 1 \ }
\end{equation}
We additionally set $\Nmu(X,X) = 0$, and $\Nmu(X,Y) = \infty$ when $X$ and $Y$ lie in different trees, meaning no sequence of matings connects them at all.

In the below figure, we observe a pedigree structure where $\Nmu = 0, 1, 2, 3, \infty$ can be observed:

\begin{figure}[H]
    \begin{floatcard}
        \definecolor{pmCB03A2E}{HTML}{B03A2E}
\begin{tikzpicture}[
  every node/.style={font=\small},
  pmObserved/.style={draw=black, fill=white, rectangle, inner sep=0.06cm, minimum width=0.38cm, minimum height=0.34cm},
  pmLatent/.style={draw=black, fill=white, ellipse, inner xsep=0.07cm, inner ysep=0.13cm, minimum width=0.38cm, minimum height=0.34cm},
  pmDirected/.style={-{Stealth[length=2mm]}, line width=0.7pt},
  pmBidirected/.style={{Stealth[length=2mm]}-{Stealth[length=2mm]}, line width=0.7pt},
  pmCopath/.style={line width=1.6pt},
  pmLabel/.style={midway, fill=white, inner sep=1pt},
  pmLabelAt/.style={fill=white, inner sep=1pt},
]
  \node[pmLatent] (g_A) at (0.000,3.400) {$g_{A}$};
  \node[pmLatent] (g_B) at (3.800,3.400) {$g_{B}$};
  \node[pmLatent] (g_C) at (7.600,3.400) {$g_{C}$};
  \node[pmLatent] (g_D) at (11.400,3.400) {$g_{D}$};
  \node[pmLatent] (g_E) at (1.900,1.800) {$g_{E}$};
  \node[pmLatent] (g_F) at (5.700,1.800) {$g_{F}$};
  \node[pmLatent] (g_G) at (9.500,1.800) {$g_{G}$};
  \node[pmObserved] (y_A) at (0.000,5.200) {$y_{A}$};
  \node[pmObserved] (y_B) at (3.800,5.200) {$y_{B}$};
  \node[pmObserved] (y_C) at (7.600,5.200) {$y_{C}$};
  \node[pmObserved] (y_D) at (11.400,5.200) {$y_{D}$};
  \node[pmObserved] (y_E) at (1.900,0.000) {$y_{E}$};
  \node[pmObserved] (y_F) at (5.700,0.000) {$y_{F}$};
  \node[pmObserved] (y_G) at (9.500,0.000) {$y_{G}$};
  \draw[pmDirected, draw=black] (g_A) -- (g_E);
  \draw[pmDirected, draw=black] (g_B) -- (g_E);
  \draw[pmDirected, draw=black] (g_B) -- (g_F);
  \draw[pmDirected, draw=black] (g_C) -- (g_F);
  \draw[pmDirected, draw=black] (g_C) -- (g_G);
  \draw[pmDirected, draw=black] (g_D) -- (g_G);
  \draw[pmDirected, draw=black] (g_A) -- (y_A);
  \draw[pmDirected, draw=black] (g_B) -- (y_B);
  \draw[pmDirected, draw=black] (g_C) -- (y_C);
  \draw[pmDirected, draw=black] (g_D) -- (y_D);
  \draw[pmDirected, draw=black] (g_E) -- (y_E);
  \draw[pmDirected, draw=black] (g_F) -- (y_F);
  \draw[pmDirected, draw=black] (g_G) -- (y_G);
  \draw[pmBidirected, draw=black] (g_A) to[loop, looseness=6, in=210, out=150] (g_A);
  \node[pmLabelAt] at (-0.855,3.400) {$V_{A}$};
  \draw[pmBidirected, draw=black] (g_B) to[loop, looseness=6, in=75, out=15] (g_B);
  \node[pmLabelAt] at (4.519,4.112) {$V_{A}$};
  \draw[pmBidirected, draw=black] (g_C) to[loop, looseness=6, in=75, out=15] (g_C);
  \node[pmLabelAt] at (8.319,4.112) {$V_{A}$};
  \draw[pmBidirected, draw=black] (g_D) to[loop, looseness=6, in=30, out=330] (g_D);
  \node[pmLabelAt] at (12.255,3.400) {$V_{A}$};
  \draw[pmBidirected, draw=black] (g_E) to[loop, looseness=11, in=120, out=60] (g_E);
  \node[pmLabelAt] at (1.900,3.145) {$\tfrac12 V_A\big(1-\rho_g\big)$};
  \draw[pmBidirected, draw=black] (g_F) to[loop, looseness=11, in=120, out=60] (g_F);
  \node[pmLabelAt] at (5.700,3.145) {$\tfrac12 V_A\big(1-\rho_g\big)$};
  \draw[pmBidirected, draw=black] (g_G) to[loop, looseness=11, in=120, out=60] (g_G);
  \node[pmLabelAt] at (9.500,3.145) {$\tfrac12 V_A\big(1-\rho_g\big)$};
  \draw[pmBidirected, draw=black] (y_A) to[loop above, looseness=6, in=120, out=60] node[pmLabel] {$V_{E}$} (y_A);
  \draw[pmBidirected, draw=black] (y_B) to[loop above, looseness=6, in=120, out=60] node[pmLabel] {$V_{E}$} (y_B);
  \draw[pmBidirected, draw=black] (y_C) to[loop above, looseness=6, in=120, out=60] node[pmLabel] {$V_{E}$} (y_C);
  \draw[pmBidirected, draw=black] (y_D) to[loop above, looseness=6, in=120, out=60] node[pmLabel] {$V_{E}$} (y_D);
  \draw[pmBidirected, draw=black] (y_E) to[loop, looseness=6, in=30, out=330] (y_E);
  \node[pmLabelAt] at (2.651,0.000) {$V_{E}$};
  \draw[pmBidirected, draw=black] (y_F) to[loop, looseness=6, in=30, out=330] (y_F);
  \node[pmLabelAt] at (6.451,0.000) {$V_{E}$};
  \draw[pmBidirected, draw=black] (y_G) to[loop, looseness=6, in=30, out=330] (y_G);
  \node[pmLabelAt] at (10.251,0.000) {$V_{E}$};
  \draw[pmCopath, draw=pmCB03A2E] (y_A) -- node[pmLabel] {$[\rho_y]$} (y_B);
  \draw[pmCopath, draw=pmCB03A2E] (y_B) -- node[pmLabel] {$[\rho_y]$} (y_C);
  \draw[pmCopath, draw=pmCB03A2E] (y_C) -- node[pmLabel] {$[\rho_y]$} (y_D);
\end{tikzpicture}

        \caption{%
            A pedigree of four parents joined by three matings, with one child per mating, drawn at the level of genetic values as in Figure~\ref{fig:g-only}.
            Reading left to right, $A$ and $B$ are the parents of $E$, $B$ and $C$ of $F$, and $C$ and $D$ of $G$.
            Each parent is $\Nmu = 0$ to themselves.
            Each parent is $\Nmu = 1$ to their own mate, e.g.\ $(A,B)$; $\Nmu = 2$ to their mate's other mate, e.g.\ $(A,C)$ and $(B,D)$; and $\Nmu = 3$ across the whole chain, $(A,D)$.
            In the offspring generation, by contrast, every pair is $\Nmu = \infty$: none of $E$, $F$, $G$ has mated, so no chain of matings connects any two of them.%
        }
        \label{fig:step-sib}
    \end{floatcard}
\end{figure}

\subsubsection{Lineal chains}
The \emph{meiosis graph} consists only of the directed edges (parent-to-offspring meioses) between nodes.
Every individual has two edges coming in, one from each parent, or none at all if they are a founder of the pedigree.
A \emph{directed path} is a sequence of distinct nodes in which each is joined by an edge to the next, always with the edge pointing forward.
Under \assumptionid{Discrete-generations} every edge steps from one generation to the next, so following edges forward always moves forward in time.
Therefore, the meiosis graph is guaranteed to be a \emph{directed acyclic graph}.

Let $t_X$ denote the generation that an individual $X$ belongs to, according to \assumptionid{Discrete-generations}.
Let $m(X)$ and $p(X)$ denote the two parents, mother and father, of an individual $X$, i.e. the two nodes with an edge into $X$.
If a directed path runs from $Y$ to $X$, then $Y$ is an \textbf{ancestor} of $X$, $X$ is a \textbf{descendant} of $Y$, and the two are \textbf{lineal relatives} to each other.

The \assumptionid{No-inbreeding} assumption from earlier effectively constrains our meiosis graph such that there is at most one directed path between any two individuals.
Since any two individuals must always share a common ancestor, this can never be true in theory.
However, in practice, maintaining this constraint for a couple generations back should be sufficient.

For lineal relatives $X$ and $Y$, the \textbf{lineal chain} $\Lin(X,Y)$ is the unique directed path between them, read as an ordered list running from $X$ to $Y$.
The \textbf{lineal distance} $\Nlin(X,Y)$ is the number of edges on that path, i.e. the number of meioses separating the two.
Under \assumptionid{Discrete-generations}, it is just the number of generations between them:
    \begin{equation}\label{eq:lineal-dist}
        \boxed{\ \Nlin(X,Y) \;=\; \big|\Lin(X,Y)\big| - 1 \;=\; \big|\, t_X - t_Y \,\big| \ }
    \end{equation}
We additionally set $\Nlin(X,X) = 0$, and $\Nlin(X,Y) = \infty$ if $X$ and $Y$ are not lineal relatives.

\subsubsection{Tracing paths between relatives}


    A \textbf{path} $\mathcal{P}(A,B)$ between probands $A$ and $B$ is a sequence of chains, alternating between lineal and mating, in which each chain begins at the node where the previous one ends.
    The first chain begins at node $A$, the last ends at $B$, and no individual is visited twice across the entire path.
    We can think of a path itself as being a \emph{chain digraph}, where the nodes are chains. and the edges connect lineal chains to mating chains, or vice-versa.
    Since mating chains contain individuals of the same generation while lineal chains contain individuals spanning multiple generations, we enforce edges to be directed such that the older chain points to the younger chain.
    The number of edges leaving a chain node is its out-degree, $N_o \in \{0,1,2\}$, and the number of edges entering it is its in-degree.
    Lineal chains will always have an out-degree of 1 and an in-degree of 1.

    Furthermore, we require that whenever a mating chain has an out-edge, the mating chain must contain both the parents of the second-oldest individual in the lineal chain it points to.
    In other words, the terminal individual node of two connected lineal and mating chains must be the ``outer'' parent.
    One can think of an out-edge from a mating chain as representing a mate pair, rather than an individual.


Below is a pedigree where the path between $A$ and $B$ is traced, which spans 2 mating chains and 3 lineal chains, followed by its equivalent digraph.

\begin{figure}[H]
    \begin{floatcard}
        \makebox[\linewidth][l]{\textbf{(A)}}
        \definecolor{pmCPath9}{HTML}{999999}
\definecolor{pmCPathR}{HTML}{B03A2E}
\definecolor{pmCPathB}{HTML}{1F77B4}
\begin{tikzpicture}[
  every node/.style={font=\small},
  pmPerson/.style={draw=black, fill=pmCPathB!12, circle, inner sep=0.04cm,
                   minimum size=0.60cm},
  pmPersonOff/.style={draw=pmCPath9, text=pmCPath9, fill=white, circle,
                      inner sep=0.04cm, minimum size=0.60cm},
  pmDirected/.style={-{Stealth[length=2mm]}, line width=0.9pt},
  pmDirectedOff/.style={-{Stealth[length=1.7mm]}, line width=0.7pt, draw=pmCPath9},
  pmCopath/.style={line width=1.6pt, draw=pmCPathR},
  pmCopathOff/.style={line width=1.6pt, draw=pmCPath9},
  pmGen/.style={font=\footnotesize, text=pmCPath9, anchor=east},
]
  \node[pmGen] at (-0.700,4.500) {$t$};
  \node[pmGen] at (-0.700,3.000) {$t+1$};
  \node[pmGen] at (-0.700,1.500) {$t+2$};
  \node[pmGen] at (-0.700,0.000) {$t+3$};
  \node[pmPerson] (I) at (5.500,4.500) {$I$};
  \node[pmPerson] (J) at (6.500,4.500) {$J$};
  \node[pmPerson] (K) at (8.000,4.500) {$K$};
  \node[pmPerson] (L) at (9.500,4.500) {$L$};
  \node[pmPerson] (A) at (0.000,3.000) {$A$};
  \node[pmPersonOff] (M) at (2.000,3.000) {$M$};
  \node[pmPersonOff] (N) at (5.000,3.000) {$N$};
  \node[pmPerson] (H) at (6.000,3.000) {$H$};
  \node[pmPerson] (B) at (8.750,3.000) {$B$};
  \node[pmPerson] (C) at (1.000,1.500) {$C$};
  \node[pmPersonOff] (O) at (3.000,1.500) {$O$};
  \node[pmPersonOff] (P) at (4.500,1.500) {$P$};
  \node[pmPerson] (G) at (5.500,1.500) {$G$};
  \node[pmPerson] (D) at (2.000,0.000) {$D$};
  \node[pmPerson] (E) at (3.500,0.000) {$E$};
  \node[pmPerson] (F) at (5.000,0.000) {$F$};
  \draw[pmDirected] (A) -- (C);
  \draw[pmDirectedOff] (M) -- (C);
  \draw[pmDirected] (C) -- (D);
  \draw[pmDirectedOff] (O) -- (D);
  \draw[pmDirected] (I) -- (H);
  \draw[pmDirectedOff] (J) -- (H);
  \draw[pmDirected] (H) -- (G);
  \draw[pmDirectedOff] (N) -- (G);
  \draw[pmDirected] (G) -- (F);
  \draw[pmDirectedOff] (P) -- (F);
  \draw[pmDirectedOff] (K) -- (B);
  \draw[pmDirected] (L) -- (B);
  \draw[pmCopathOff] (A) -- (M);
  \draw[pmCopathOff] (C) -- (O);
  \draw[pmCopathOff] (G) -- (P);
  \draw[pmCopathOff] (H) -- (N);
  \draw[pmCopath] (I) -- (J);
  \draw[pmCopath] (J) -- (K);
  \draw[pmCopath] (K) -- (L);
  \draw[pmCopath] (D) -- (E);
  \draw[pmCopath] (E) -- (F);
\end{tikzpicture}


        \vspace{1.0em}
        \makebox[\linewidth][l]{\textbf{(B)}}
        \definecolor{pmCPathB}{HTML}{1F77B4}
\definecolor{pmCPath9}{HTML}{999999}
\begin{tikzpicture}[
  every node/.style={font=\small},
  cdMating/.style={draw=black, fill=pmCPathB!14, rectangle, rounded corners=1pt,
                   inner sep=0.10cm, minimum height=0.62cm},
  cdLineal/.style={draw=black, fill=white, ellipse,
                   inner xsep=0.06cm, inner ysep=0.10cm, minimum height=0.62cm},
  cdProband/.style={draw=black, fill=pmCPathB!14, circle, inner sep=0.04cm,
                    minimum size=0.60cm},
  cdEdge/.style={-{Stealth[length=2mm]}, line width=0.9pt},
  cdNote/.style={font=\scriptsize, text=pmCPath9, align=center, anchor=north},
]
  \node[cdProband] (pA) at (0.000,0.000) {$A$};
  \node[cdLineal] (L1) at (2.150,0.000) {$\Lin(A,D)$};
  \node[cdNote] at (2.150,-0.460) {$\Nlin = 2$};
  \node[cdMating] (M1) at (4.300,0.000) {$\Mu(D,F)$};
  \node[cdNote] at (4.300,-0.460) {$\Nmu = 2$\\$N_o = 0$};
  \node[cdLineal] (L2) at (6.450,0.000) {$\Lin(F,I)$};
  \node[cdNote] at (6.450,-0.460) {$\Nlin = 3$};
  \node[cdMating] (M2) at (8.600,0.000) {$\Mu(I,L)$};
  \node[cdNote] at (8.600,-0.460) {$\Nmu = 3$\\$N_o = 2$};
  \node[cdLineal] (L3) at (10.750,0.000) {$\Lin(L,B)$};
  \node[cdNote] at (10.750,-0.460) {$\Nlin = 1$};
  \node[cdProband] (pB) at (12.900,0.000) {$B$};
  \draw[cdEdge] (pA) -- (L1);
  \draw[cdEdge] (L1) -- (M1);
  \draw[cdEdge] (L2) -- (M1);
  \draw[cdEdge] (M2) -- (L2);
  \draw[cdEdge] (M2) -- (L3);
  \draw[cdEdge] (L3) -- (pB);
\end{tikzpicture}

        \caption{%
            (A) A pedigree shown as a graph containing a highlighted path connecting $A$ and $B$.
            Meioses are denoted by black directed edges while matings are denoted by red undirected edges.
            (B) The chain digraph equivalent of the highlighted path:
            $\mathcal{P}(A,B) = \big(\Lin(A,D),\ \Mu(D,F),\ \Lin(F,I),\ \Mu(I,L),\ \Lin(L,B)\big)$.
            Edges are oriented to always point from the older chain to the younger chain.
            Properties of each chain are listed.
            Note that proband $A$ acts as a source while proband $B$ acts as a sink.
        }
        \label{fig:path-example}
    \end{floatcard}
\end{figure}

\subsubsection{Covariance along a single chain}
Before assembling a whole path, we derive the covariance carried by a single chain of each type.

Two individuals $X$ and $Y$ on a mating chain are joined by exactly $\Nmu(X,Y)$ matings.
Each mating crossed contributes a factor of $\mu$ by \eqref{eq:leg-rule}, and each of the $\Nmu(X,Y) - 1$ individuals passed through in between contributes their own $\Var[y] = V_Y$, as does each end:
\begin{equation}\label{eq:chain-y}
\begin{gathered}
    \Cov[\,y_X,\, y_Y\,] \;=\; V_Y \cdot \mu \cdot
    \underbrace{V_Y \cdot \mu \;\cdots\; V_Y \cdot \mu}_{\Nmu(X,Y) - 1 \text{ times}} \cdot V_Y
    \;=\; \mu^{\Nmu(X,Y)}\, V_Y^{\,\Nmu(X,Y) + 1} \\[8pt]
    \boxed{\ \Cov[\,y_X,\, y_Y\,] \;=\; \rho_y^{\,\Nmu(X,Y)}\, V_Y \ }
\end{gathered}
\end{equation}
This holds even when $X = Y$.
Reading the two ends at the genetic level instead exchanges the $V_Y$ at each end for a $V_A$, which can also be simplified.
Unlike before, this general formula does not hold for $X = Y$, since an individual's genetic value has covariance $V_A$ with itself rather than $h^2 V_A$:
\begin{equation}\label{eq:chain-g}
    \boxed{\
    \Cov[\,g_X,\, g_Y\,] \;=\;
    \begin{cases}
        V_A, & \Nmu(X,Y) = 0 \\[6pt]
        \rho_g\, \rho_y^{\,\Nmu(X,Y) - 1}\, V_A, & \Nmu(X,Y) \geq 1
    \end{cases} \ }
\end{equation}

For a lineal chain, the covariance is an extension of the parent--offspring column of Table~\ref{tab:g-only-cov}, contributing one transmission coefficient per meiosis:
\begin{equation}\label{eq:lineal-g}
\begin{gathered}
    \Cov[\,y_X,\, g_{Y}\,] \;=\; \Cov[\,g_X,\, g_{Y}\,] \;=\; V_A \cdot
    \underbrace{\frac{1+\rho_g}{2} \cdot \frac{1+\rho_g}{2} \;\cdots\; \frac{1+\rho_g}{2}}_{\Nlin(X,Y) \text{ meioses}} \\[8pt]
    \boxed{\ \Cov[\,y_X,\, g_{Y}\,] \;=\; \Cov[\,g_X,\, g_{Y}\,] \;=\; V_A
    \left(\frac{1+\rho_g}{2}\right)^{\!\Nlin(X,Y)} \ }
\end{gathered}
\end{equation}
If we read the ancestor $Y$ at their phenotypic value, we have to account for the fact that the environment of $Y$ informs the genetic value that the mate of $Y$ transmits to their offspring $X$.
As we saw in Table~\ref{tab:g-only-cov}, this is achieved by modifying the coefficient of the first meiosis step:
\begin{equation}\label{eq:lineal-y}
\begin{gathered}
    \Cov[\,y_X,\, y_{Y}\,] \;=\; \Cov[\,g_X,\, y_{Y}\,] \;=\; V_A \cdot
    \underbrace{\frac{1+\rho_y}{2}}_{\text{first meiosis}} \cdot
    \underbrace{\frac{1+\rho_g}{2} \;\cdots\; \frac{1+\rho_g}{2}}_{\Nlin(X,Y) - 1 \text{ meioses}} \\[8pt]
    \boxed{\ \Cov[\,y_X,\, y_{Y}\,] \;=\; \Cov[\,g_X,\, y_{Y}\,] \;=\; V_A
    \left(\frac{1+\rho_y}{2}\right)
    \left(\frac{1+\rho_g}{2}\right)^{\!\Nlin(X,Y) - 1} \ }
\end{gathered}
\end{equation}

\subsubsection{Covariance between probands on a path}\label{sec:am-probands}

We now want to express the product of the covariance path coefficients implied by our traced pedigree path between $A$ and $B$ as a function of the properies of the lineal and mating chains contained in it.
Furthermore, this path coefficient product should then roughly equal the covariance between the two probands.
Enforcing this requires the following assumption and approximation:

\begin{assumption}{Outbred-mating-chain}{A path's lineal chains account for all of a pair's relatedness}
    The expected proportion of the genome that two individuals on a path share IBD, their realized relatedness $\pi_{ij}$ of \eqref{eq:pi-def}, given the full pedigree should be equal to that which is explained solely by meiotic inheritance from common ancestor(s) lying on those individuals' lineal chains.

    More specifically, if $X$ and $Y$ are lineal relatives, then all of the IBD segments between them must have been transmitted down a lineal chain that connects them.
    If $X$ and $Y$ are not lineal relatives but lie on different lineal chains present in the path such that mating chain $\mathcal{M}_{XY}$ has an out-edge pointing to both lineal chains, then any IBD segments between $X$ and $Y$ must have been inherited from a common ancestor present in $\mathcal{M}_{XY}$.
\end{assumption}

\begin{approximation}{Dominant-path}{A single path dominates the covariance between two individuals}
    Since any two individuals must share a common ancestor when looking up the pedigree far enough, there is never a unique path that connects the two.
    The true covariance of the pair will be the sum of the product of path coefficients across all paths.

    We approximate the covariance between the pair, given the full pedigree, to be the single largest of those products, and we call the path carrying it the \emph{dominant path}.
    Whenever $\rho_y \geq 0$ every product we discard is itself non-negative, so the approximation is a guaranteed lower bound.
\justification{
    Because the path coefficient of additional steps in a pedigree are between 0 and 1, the dominant path tends to be the shortest path by steps.
    Pedigrees that break \assumptionid{No-inbreeding} or \assumptionid{Outbred-mating-chain} only add further valid paths, and those are guaranteed to contain more steps than the dominant one.
    Therefore, its contribution to the total variance would be much smaller than that of the dominant path.
    }
\end{approximation}

Let $i$ index the lineal chains in a path and $j$ index the mating chains in a path.
We can then write the following formula for the covariance in genetic values of any two individuals:

\begin{equation}\label{eq:path-product}
    \boxed{\ \Cov[\,g_A,\, g_B\,] \;=\; V_A\,
    \left( \frac{1+\rho_g}{2} \right)^{\textstyle \sum_i N_{\Lin_i} \;-\; \sum_j N_{o_j}}
    \ \prod_j\, w_j \ }
\end{equation}

We briefly explain some of the terms in the expression.
$V_A$ comes from the variance of $g$ itself.
$\tfrac12(1+\rho_g)$ comes from the parent--offspring coefficient of Table~\ref{tab:g-only-cov} written as a transmission rate.
Its exponent counts the number of ascents/descents across all lineal paths, $\sum_i N_{\Lin_i} $, but removes the double-counting that occurs when lineal chains ascend into mating chains, $N_{o_j}$.
$\prod_j\, w_j$ comes from the crossing of mating paths and the lineal chains adjacent to them, explained below.

The weight $w_j$ of a mating chain is the implied genetic covariance $c$ averaged across the possible ``entry points'' into the mating chain from the adjacent lineal chains.

\begin{equation}\label{eq:chain-weight}
\begin{gathered}
    c(\Nmu) \;=\;
    \begin{cases}
        1 & \Nmu = 0 \\[3pt]
        \rho_g\, \rho_y^{\,\Nmu-1} & \Nmu \geq 1
    \end{cases}
    \\[12pt]
    \boxed{\ w_j \;=\; \frac{1}{2^{N_{o_j}}}
    \sum_{n \,=\, 0}^{N_{o_j}}
    \binom{N_{o_j}}{n}\;
    c\Big( \big|\, N_{\Mu_j} \,-\, n \,\big| \Big) \ }
\end{gathered}
\end{equation}


Here $c(\Nmu)$ is the correlation in genetic values between two individuals at mating distance $\Nmu$, derived above in \eqref{eq:chain-g}.
The averaging is over the mating chain's two ends, and the two kinds of end differ in how many individuals the path enters at.
Along an in-edge the path reaches the mating chain at a single individual.
Along an out-edge it reaches a mate pair, since we require the mating chain to carry both parents of the second-oldest individual of the lineal chain it points to.
Taking only the outer parent of that pair would understate the product: an offspring's genetic value comes half from each parent, and the inner parent sits one mating nearer the far end of the chain, so their half is worth the more of the two.

In real data, dominant paths' mating chains tend to be sources, i.e. $N_o=2$.
In such cases, the mating chain is guaranteed to contain both parents of the second-oldest individual of each adjacent lineal chain.
We categorize the relationship type of these individuals as \emph{generalized siblings}.
We define $N_\mathcal{S} \in \{0,1,2\}$ to be the number of shared parents between generalized siblings.
$N_\mathcal{S} = 2$ indicates \emph{full-siblings},
$N_\mathcal{S} = 1$ indicates \emph{half-siblings}, and 
$N_\mathcal{S} = 0$ indicates \emph{generalized step-siblings},
for which the $N_\Mu = 3$ case is dubbed just as \emph{step-siblings}.
Under this parameterization, where the first two cases require a source mating chain and the third holds at any $N_{o_j}$, we can write:

\begin{equation}\label{eq:chain-weight-cases}
    w_j \;=\;
    \begin{cases}
        \dfrac{1+\rho_g}{2}
            & N_{\mathcal{S}_j} = 2, \quad N_{o_j} = 2 \\[14pt]
        \dfrac{1 + 2\rho_g + \rho_g\rho_y}{4}
            & N_{\mathcal{S}_j} = 1, \quad N_{o_j} = 2 \\[14pt]
        \rho_g\; \rho_y^{\,N_{\Mu_j} - 1 - N_{o_j}}
        \left( \dfrac{1+\rho_y}{2} \right)^{\!N_{o_j}}
            & N_{\mathcal{S}_j} = 0, \quad \text{any } N_{o_j}
    \end{cases}
\end{equation}

The weight for the half-sibling case ($N_\mathcal{S} = 1$) is Nagylaki's half-sibling multiplier \citep{Nagylaki1978}.
The third case is the \eqref{eq:chain-weight} polynomial expanded under the condition that $N_{\Mu_j} \geq N_{o_j} + 1$, which is exactly what $N_{\mathcal{S}_j} = 0$ encodes, so that no term of \eqref{eq:chain-weight} reaches $c(0)$.

Lastly, we have to include a correction factor for each proband for which we are interested in the covariance with the phenotype, as opposed to the genetic value, based on the position of the proband in the path.

\begin{equation}\label{eq:proband-correction}
    \boxed{\ C_X \;=\;
    \begin{cases}
        1
            & \text{$X$ is a sink} \\[10pt]
        \dfrac{1+\rho_y}{1+\rho_g}
            & \text{$X$ is a source} \\[10pt]
        \dfrac{\rho_y}{\rho_g} \;=\; \dfrac{1}{h^2}
            & \text{$X$ is in a terminal mating chain}
    \end{cases} \ }
\end{equation}

The three cases are the consequences of applying the rule that switching from a proband's genetic value to their phenotypic value only requires replacing $\rho_g$ by $\rho_y$ in that proband's own mating and meiosis edges.
A proband that is a sink is the youngest individual on its lineal chain, while a proband that is a source is the oldest individual on its lineal chain.
A proband that is part of a terminal mating chain has no lineal chain of its own at all.
We can therefore express the covariance between the phenotypes of $A$ and $B$ as:

\begin{equation}\label{eq:path-product-pheno}
    \boxed{\ \Cov[\,y_A,\, y_B\,] \;=\; C_A\; C_B\; \Cov[\,g_A,\, g_B\,] \ }
\end{equation}

\subsubsection{Example}
Below, we apply our rules to the example shown in Figure~\ref{fig:path-example}.
In its chain digraph, $\Mu(D,F)$ is a sink and so has $N_o = 0$, while $\Mu(I,L)$ is a source and so has $N_o = 2$.
The two chains are $N_\Mu = 2$ and $N_\Mu = 3$ long, and neither connects siblings, so $N_\mathcal{S} = 0$ for both.
Reading the two counts off the path and then the weight of each mating chain:

\begin{equation*}
\begin{aligned}
    \sum_i N_{\Lin_i} \;&=\; \Nlin(A,D) + \Nlin(F,I) + \Nlin(L,B) \;=\; 2 + 3 + 1 \;=\; 6 \\[8pt]
    \sum_j N_{o_j} \;&=\; 0 \;+\; 2 \;=\; 2 \\[8pt]
    w_{\Mu(D,F)} \;&=\; \rho_g\, \rho_y^{\,2-1-0} \;=\; \rho_g \rho_y \\[8pt]
    w_{\Mu(I,L)} \;&=\; \rho_g\, \rho_y^{\,3-1-2} \left(\frac{1+\rho_y}{2}\right)^{\!2}
    \;=\; \rho_g \left(\frac{1+\rho_y}{2}\right)^{\!2}
\end{aligned}
\end{equation*}

Substituting these into \eqref{eq:path-product}:

\begin{equation*}
\begin{aligned}
    \Cov[\,g_A,\, g_B\,]
    \;&=\; V_A \left(\frac{1+\rho_g}{2}\right)^{\!6-2}
    \cdot \rho_g\rho_y
    \cdot \rho_g \left(\frac{1+\rho_y}{2}\right)^{\!2} \\[8pt]
    \;&=\; V_A\; \rho_g^2\, \rho_y
    \left(\frac{1+\rho_g}{2}\right)^{\!4}
    \left(\frac{1+\rho_y}{2}\right)^{\!2}
\end{aligned}
\end{equation*}

The two probands take different corrections.
$A$ is the oldest individual on $\Lin(A,D)$, so $C_A = (1+\rho_y)/(1+\rho_g)$, while $B$ is the youngest on $\Lin(L,B)$ and takes $C_B = 1$:
\begin{equation*}
    \Cov[\,y_A,\, y_B\,] \;=\; \frac{1+\rho_y}{1+\rho_g}\; \Cov[\,g_A,\, g_B\,]
    \;=\; V_A\; \rho_g^2\, \rho_y
    \left(\frac{1+\rho_g}{2}\right)^{\!3}
    \left(\frac{1+\rho_y}{2}\right)^{\!3}
\end{equation*}

\subsubsection{Table of covariances for common relationship types}

\begin{table}[H]
    \begin{floatcard}
    \renewcommand{\arraystretch}{2.1}
    \setlength{\tabcolsep}{4pt}
    \scriptsize
    \begin{tabular}{l|cc|c|c|cc}
        relationship
            & $\displaystyle\sum_i N_{\Lin_i}$
            & $\displaystyle\sum_j N_{o_j}$
            & $\prod_j w_j$
            & $\Cov[\,g_A,\, g_B\,]$
            & $C_A$ & $C_B$ \\
        \hline
        self
            & $0$ & $0$ & $1$ & $V_A$
            & --- & --- \\
        parent--parent
            & $0$ & $0$ & $\rho_g$ & $\rho_g V_A$
            & $\tfrac{\rho_y}{\rho_g}$ & $\tfrac{\rho_y}{\rho_g}$ \\
        parent--offspring
            & $1$ & $0$ & $1$
            & $\big(\tfrac{1+\rho_g}{2}\big) V_A$
            & $\tfrac{1+\rho_y}{1+\rho_g}$ & $1$ \\
        grandparent--grandchild
            & $2$ & $0$ & $1$
            & $\big(\tfrac{1+\rho_g}{2}\big)^{2} V_A$
            & $\tfrac{1+\rho_y}{1+\rho_g}$ & $1$ \\
        full siblings
            & $2$ & $2$ & $\tfrac{1+\rho_g}{2}$
            & $\big(\tfrac{1+\rho_g}{2}\big) V_A$
            & $1$ & $1$ \\
        avuncular
            & $3$ & $2$ & $\tfrac{1+\rho_g}{2}$
            & $\big(\tfrac{1+\rho_g}{2}\big)^{2} V_A$
            & $1$ & $1$ \\
        first cousins
            & $4$ & $2$ & $\tfrac{1+\rho_g}{2}$
            & $\big(\tfrac{1+\rho_g}{2}\big)^{3} V_A$
            & $1$ & $1$ \\
        half siblings
            & $2$ & $2$ & $\tfrac{1 + 2\rho_g + \rho_g\rho_y}{4}$
            & $\tfrac{1 + 2\rho_g + \rho_g\rho_y}{4} V_A$
            & $1$ & $1$ \\
        step siblings
            & $2$ & $2$ & $\rho_g\big(\tfrac{1+\rho_y}{2}\big)^{2}$
            & $\rho_g\big(\tfrac{1+\rho_y}{2}\big)^{2} V_A$
            & $1$ & $1$ \\
        step-step siblings
            & $2$ & $2$ & $\rho_g\rho_y\big(\tfrac{1+\rho_y}{2}\big)^{2}$
            & $\rho_g\rho_y\big(\tfrac{1+\rho_y}{2}\big)^{2} V_A$
            & $1$ & $1$ \\
        parents of a married couple
            & $2$ & $0$ & $\rho_g$
            & $\rho_g\big(\tfrac{1+\rho_g}{2}\big)^{2} V_A$
            & $\tfrac{1+\rho_y}{1+\rho_g}$ & $\tfrac{1+\rho_y}{1+\rho_g}$ \\
    \end{tabular}
    \caption{%
        Common relationship classes, each assembled from \eqref{eq:path-product}.
        The mating chain of the sibling rows runs between the two outer parents, at $\Nmu = 1, 2, 3, 4$ for full, half, step and step-step siblings respectively; the classes are therefore the same path with one further mating crossed each time.
        Two coincidences are worth noting, because they arrive from different rows of \eqref{eq:chain-weight-cases}: full siblings match parent--offspring, and avuncular matches grandparent--grandchild.
        The \emph{self} row is degenerate --- there is no path, so no proband correction applies and $\Cov[\,y_X,\, y_X\,]$ is simply $V_Y$.%
    }
    \label{tab:relationship-classes}
    \end{floatcard}
\end{table}

\comment{For agents, the below is to not be edited yet. The user will explicitly say when to edit below this point. The text below this was written previously but is not guaranteed to be correct. This region ends at the matching comment just above the references.}

\subsection{Level 0 --- unchanged}
Conditioning on both individuals' causal genotypes and the true effect sizes leaves nothing to average over, so \eqref{eq:level0} holds exactly as written.
Level 0 is a statement about the genotype-to-phenotype map alone, and no mating rule can disturb it.
That it survives every architecture untouched is a feature of where the tower starts, not a coincidence.

\subsection{Level 1 --- the off-diagonal no longer collapses}\label{sec:am-level1}
Section 1 moved from Level 0 to Level 1 by averaging over the effects and observing that the off-diagonal $k \neq k'$ terms vanish.
Two separate clauses of \assumptionid{Random-effects} were needed for that: the effects are marginally uncorrelated, $\E[\beta_k\beta_{k'}]=0$, and their cross-products are independent of the genotypes, $\E[\beta_k\beta_{k'}\mid\bX] = \E[\beta_k\beta_{k'}]$.

The first clause survives assortative mating.
Assortment does not reach back and change the effect sizes; it changes which genotypes are common.

The second clause does not survive, and it is worth seeing why the failure is unavoidable rather than incidental.
Under assortment, the correlations that build up between variants are generated \emph{by} the effects: it is precisely the trait-increasing alleles that end up travelling together.
The observed pattern of correlation in $\bX$ is therefore informative about which variants have effects and with what signs, so conditioning on $\bX$ shifts our belief about $\beta_k\beta_{k'}$.
Formally, $\E[\beta_k \beta_{k'} \mid \bX]$ is positive for pairs of variants that are positively correlated in the sample, and the off-diagonal sum no longer disappears:
\begin{equation}\label{eq:am-level1}
    \E[\,g_i g_j \mid \bx_i, \bx_j\,]
    \;=\; \underbrace{\sum_k \E[\beta_k^2]\, x_{ik} x_{jk}}_{\textstyle \VAo A_{ij}}
    \;+\; \underbrace{\sum_{k}\sum_{k'\neq k} \E[\beta_k \beta_{k'}\mid \bX]\, x_{ik} x_{jk'}}_{\textstyle \text{disequilibrium term}} .
\end{equation}
The first term is the Level-1 result of Section 1, with the additive variance replaced by the \emph{genic} variance $\VAo$ of \eqref{eq:am-VAsplit} --- the genotype matrix on its own can only recover the per-variant contributions.
The second term is new, and it is what the remainder of the disequilibrium $D$ lives in.

Note the shape of \eqref{eq:am-level1}: two relatedness-like quantities added together, not one quantity rescaled.
This is the form that methods for unrelated individuals take, where the second term is estimated with a second relatedness matrix alongside the GRM \citep{Zhang2025NatGenet}.
We do not develop it here.

\subsection{The tower re-orders}
At this point Section 1 descended to realized relatedness and then to genetic degree.
That route is no longer available, and the reason is the screening property named at the end of Section 1.

Under random mating, $\pi_{ij}$ screened off $d_{ij}$: once we knew what fraction of the genome a pair actually shared, their pedigree told us nothing more.
Under assortative mating this is false.
The disequilibrium contribution to a pair's covariance is carried by the correlations their \emph{ancestors} had, and how many rounds of assortment separate them --- information that lives in the pedigree and is not recoverable from how much genome they happen to share.
Two pairs with identical $\pi_{ij}$ but different pedigrees have different expected covariances.

The four information sets of Section 1 are therefore no longer a chain.
Level 2 (realized relatedness) and Level 3 (pedigree) are now two \emph{parallel} coarsenings of a common refinement --- knowing both --- and neither is a function of the other.
We keep the level numbering, since each level still refers to the same information, but the order of derivation reverses: the pedigree is now the easy case and must be done first, the joint case follows by refinement, and realized relatedness alone is the hard case and comes last.

\subsection{Level 3 --- pedigree}\label{sec:am-level3}
\subsubsection*{One more generation}
Everything at this level follows from a single recursion.

\textbf{Lemma.} \textit{Let $C$ be the child of $P$ and $P'$, and let $X$ be any random variable such that (i) $\Cov[X, e_P]=0$, and (ii) all association between $X$ and $P'$ runs through $P$'s phenotype.
Then}
\begin{equation}\label{eq:am-lemma}
    \Cov[\,X, g_C\,] \;=\; \frac{1+\rho_g}{2}\;\Cov[\,X, g_P\,] .
\end{equation}

\textit{Proof.} By \eqref{eq:am-transmission}, $\Cov[X,g_C] = \tfrac12\Cov[X,g_P] + \tfrac12\Cov[X,g_{P'}]$, the segregation term dropping out.
For the second piece, condition on $y_P$: by (ii), $\Cov[X, g_{P'}] = \Cov\big[X,\, \E[g_{P'}\mid y_P]\big]$, and the regression of a partner's genetic value on one's own phenotype has coefficient $\Cov[y_P, g_{P'}]/\VYeq = \rho_y \VAeq/\VYeq = \rho_g$ by \eqref{eq:am-rhog}.
Hence $\Cov[X,g_{P'}] = \rho_g \Cov[X, y_P] = \rho_g\Cov[X,g_P]$, the last step using (i). $\square$

Condition (ii) is a restriction on the pedigree, and it has a familiar name.

\begin{assumption}{Monogamy}{Each individual has children with one partner}
Any two relatives are connected by a single line of descent, so a chain of association between them traverses at most one assortment link that is not on that line.
\bundles{a pedigree with no half-sibships; the leading real-world violation is an individual having children with more than one partner.}
\end{assumption}

\subsubsection*{Collateral relatives}
For two \emph{collateral} relatives --- individuals related through a common ancestor, neither descended from the other --- Section~\ref{sec:am-env} showed that no environmental cross terms arise, so \eqref{eq:reduce} still applies and we need only the genetic covariance.
The base case is a pair of full siblings, children of $m$ and $f$.
Expanding both siblings with \eqref{eq:am-transmission},
\begin{align*}
    \Cov[\,g_i, g_j\,]
    &= \tfrac14\Big( \Var[g_m] + \Var[g_f] + 2\Cov[g_m,g_f] \Big) \\
    &= \tfrac14\Big( \VAeq + \VAeq + 2\rho_g \VAeq \Big)
    \;=\; \VAeq\,\frac{1+\rho_g}{2} ,
\end{align*}
using that the two siblings' transmissions from a given parent are conditionally independent given that parent, exactly as in Section 1.
Applying \eqref{eq:am-lemma} once for each further generation of separation gives, for a collateral pair of degree $d_{ij}$,
\begin{equation}\label{eq:am-level3-col}
    \boxed{\ \E[\,y_i y_j \mid d_{ij}\,] \;=\; \VAeq \left(\frac{1+\rho_g}{2}\right)^{d_{ij}} \ }
\end{equation}
in agreement with \citet{Sunde2024NatComm}.
Setting $\rho_g=0$ recovers $\VAo 2^{-d_{ij}}$, the Level-3 result of Section 1.

\subsubsection*{Lineal relatives}
For a pair in which $a$ is an ancestor of $b$, the environmental cross term of Section~\ref{sec:am-env} contributes and the base case changes.
For a parent and their offspring it is easiest to work with $\Cov[y_m, g_o]$ directly, since $e_o$ is a fresh draw:
\begin{align*}
    \Cov[\,y_m, y_o\,] = \Cov[\,y_m, g_o\,]
    &= \tfrac12\Big( \Cov[y_m,g_m] + \Cov[y_m,g_f] \Big) \\
    &= \tfrac12\Big( \VAeq + \rho_y \VAeq \Big)
    \;=\; \VAeq\,\frac{1+\rho_y}{2} .
\end{align*}
The second term is the whole story: the focal parent's \emph{phenotype} predicts the other parent's genetic value with regression coefficient $\rho_y \VAeq/\VYeq$, and half of those genes reach the child.
Iterating \eqref{eq:am-lemma} for the remaining generations gives
\begin{equation}\label{eq:am-level3-lin}
    \boxed{\ \E[\,y_a y_b \mid d_{ab}\,] \;=\; \VAeq \left(\frac{1+\rho_y}{2}\right)\left(\frac{1+\rho_g}{2}\right)^{d_{ab}-1} \ }
\end{equation}

Comparing \eqref{eq:am-level3-col} and \eqref{eq:am-level3-lin}, the two differ in one factor, $\rho_y$ in place of $\rho_g$.
Since $\rho_g = \rho_y \hsqeq < \rho_y$, the lineal covariance is always the larger.
This is a genuine and somewhat surprising consequence: under random mating a parent--offspring pair and a full-sibling pair have identical expected covariance $V_A/2$, whereas under assortative mating the parent--offspring covariance exceeds the sibling covariance by
$$ \VAeq\,\frac{\rho_y - \rho_g}{2} \;=\; \frac{\rho_y \VAeq \left(1 - \hsqeq\right)}{2} . $$
In the worked example this is a $15\%$ excess over the sibling value.
The asymmetry is purely environmental in origin --- it is the parent's own environment, correlated with the other parent's genes --- and a model that tracked only genetic covariances would miss it entirely.

\subsubsection*{Degree is no longer sufficient}
Equations \eqref{eq:am-level3-col} and \eqref{eq:am-level3-lin} disagree at the same $d$, which already shows that the degree alone does not determine the answer.
Half-siblings make the point again, from the other direction.
They are second-degree relatives, but they violate \assumptionid{Monogamy}, and the two outer parents $P_1$ and $P_2$ --- both of whom assorted on the shared parent $P$ --- are themselves correlated:
$$ \Cov[\,g_{P_1}, g_{P_2}\,] \;=\; \Var[y_P] \left(\frac{\rho_y \VAeq}{\VYeq}\right)^{\!2} \;=\; \rho_g^2\, \VYeq , $$
which gives
$$ \E[\,y_iy_j \mid \text{half-sibs}\,] \;=\; \tfrac14\Big( \VAeq (1+2\rho_g) + \rho_g^2\, \VYeq \Big) . $$
This exceeds the collateral value $\tfrac14 \VAeq (1+\rho_g)^2$ by $\tfrac14 \rho_g^2 V_E$.
So three different second-degree relationships --- grandparent--grandchild, avuncular, and half-sibling --- have three different expected covariances.

\assumptionid{Genetic-degree} must therefore be strengthened: the conditioning variable at Level 3 is not the degree but the \textbf{pedigree relationship}, of which the degree is only a summary.
The same calculation also shows that individuals with no common ancestor at all can be correlated: two people who both had children with the same person have expected phenotypic cross-product $\rho_y^2\, \VYeq$, and partners themselves $\rho_y \VYeq$.
Under random mating both are zero.

\subsection{Levels 2 and 3 together --- pedigree and realized relatedness}\label{sec:am-level23}
Knowing the pedigree fixes the \emph{expected} sharing, $\bar\pi_{d} = 2^{-d}$, but two pairs with the same relationship still differ in how much genome they actually share.
Section 1 handled that variation at Level 2; here we ask what it buys once the pedigree is already known.
Following the natural parameterisation, we write the realized sharing as its pedigree expectation plus a deviation,
$$ \pi_{ij} \;=\; \bar\pi_{d_{ij}} + \delta_{ij} , \qquad \E[\delta_{ij}\mid d_{ij}] = 0 . $$

The two terms of \eqref{eq:am-level1} respond to $\delta_{ij}$ very differently, and that is the whole content of this level.
The genic term is a sum of \emph{same-variant} products, and by \approximationid{Same-locus-AM} the argument of Section 1 goes through unchanged: its conditional expectation is $\VAo \pi_{ij}$, tracking the realized sharing exactly.
The disequilibrium term is a sum of \emph{cross-variant} products, and it behaves quite differently.

\begin{approximation}{GPD-pedigree}{The disequilibrium contribution depends only on the pedigree}
The cross-variant part of a pair's covariance is a function of their pedigree relationship alone, and does not depend on the realized sharing $\pi_{ij}$.
\justification{A cross-variant term $\E[x_{ik}x_{jk'}]$ with $k \neq k'$ is determined by \emph{which} ancestral genomes the two alleles descend from --- fixed by the pedigree --- and not by whether $i$ and $j$ happen to be IBD at either locus, since the two alleles are at different loci.
A residual dependence does exist, because a pair sharing more genome is also more likely to be IBD at $k$ \emph{and} $k'$ jointly, but for unlinked variants this is a second-order effect: the correction is of order $\rho_g \delta_{ij}$, being the product of the disequilibrium share \eqref{eq:am-Dshare} and the sharing deviation.
Deriving it exactly is left open.}
\end{approximation}

Putting the two together, the expected cross-product is the pedigree value plus a term linear in the sharing deviation.
The coefficient on that term is fixed by requiring consistency with Level 3: averaging over $\delta_{ij}$ must return the Level-3 value for that relationship.
Hence, writing $C_{ij}(\text{pedigree})$ for the Level-3 value of the pair's relationship,
\begin{equation}\label{eq:am-level23}
    \boxed{\ \E[\,y_iy_j \mid \text{pedigree},\, \pi_{ij}\,] \;=\; \underbrace{C_{ij}(\text{pedigree})}_{\text{from \eqref{eq:am-level3-col} or \eqref{eq:am-level3-lin}}} \;+\; \VAo \left(\pi_{ij} - \bar\pi_{d_{ij}}\right) \ }
\end{equation}

The structure of \eqref{eq:am-level23} is worth dwelling on, because it separates two things that Section 1 could not tell apart.
The \emph{baseline} is set by the pedigree and is inflated by assortment.
The \emph{response to realized sharing} is not inflated at all: its coefficient is $\VAo$, the base-population --- equivalently genic --- additive variance, because only the same-variant part of the covariance responds to how much genome a pair actually shares.
Under random mating the two coincide, $\VAo = \VAeq$, and \eqref{eq:am-level23} collapses back to $V_A \pi_{ij}$.

This suggests a way to see assortative mating directly, which we note without pursuing.
Within a single relationship class, regressing pairs' phenotypic cross-products on their realized relatedness gives a slope that estimates $\VAo$ and an intercept that carries the assortment-inflated baseline.
The two would agree under random mating, so the gap between them is a signature of assortment, obtainable without ever comparing across relationship classes.
Whether the residual $O(\rho_g \delta_{ij})$ term of \approximationid{GPD-pedigree} is small enough for this to be clean is the open question flagged in that box.

\subsection{Level 2 --- realized relatedness alone}\label{sec:am-level2}
The last case is the one Section 1 found easiest and assortative mating makes hardest: we observe how much genome a pair shares, but know nothing about their pedigree.
Because the pedigree is exactly what carries the disequilibrium contribution, we cannot evaluate \eqref{eq:am-level23} directly; we have to infer a pedigree from the sharing.

The natural device is to invert the degree-to-relatedness map of Section 1.
Since $\bar\pi_d = 2^{-d}$, a pair sharing a fraction $\pi_{ij}$ of their genome is, on average, as related as a relative of degree $-\log_2 \pi_{ij}$.

\begin{approximation}{Effective-degree}{Realized sharing stands in for a degree}
A pair's realized relatedness is treated as though it arose from a collateral relationship of \emph{effective degree}
$$ d^{\,\mathrm{eff}}_{ij} \;:=\; -\log_2 \pi_{ij} , $$
which is generally not an integer.
\justification{The substitution is exact whenever $\pi_{ij}$ equals its pedigree expectation for some collateral relationship, i.e.\ on the grid $\pi \in \{2^{-1}, 2^{-2}, \dots\}$; off the grid it interpolates.
The interpolation is benign for distant pairs, where many relationships of similar degree contribute and the sharing distribution is smooth, and poor for close pairs, where a handful of discrete relationships dominate and, as \S\ref{sec:am-level3} showed, relationships of the same degree can differ.
It is therefore an approximation for the distant tail, not a replacement for \eqref{eq:am-level23}.}
\end{approximation}

Substituting $d^{\,\mathrm{eff}}_{ij}$ into \eqref{eq:am-level3-col} and simplifying,
$$ \left(\frac{1+\rho_g}{2}\right)^{-\log_2 \pi_{ij}}
   = 2^{\log_2 \pi_{ij}} \cdot (1+\rho_g)^{-\log_2 \pi_{ij}}
   = \pi_{ij}\cdot \pi_{ij}^{-\log_2 (1+\rho_g)} , $$
so that the expected cross-product is a \emph{power law} in the realized relatedness:
\begin{equation}\label{eq:am-level2}
    \boxed{\ \E[\,y_i y_j \mid \pi_{ij}\,] \;=\; \VAeq\; \pi_{ij}^{\,1-\log_2(1+\rho_g)} \ }
\end{equation}

Three things are worth reading off \eqref{eq:am-level2}.

First, it contains Section 1 as a special case.
With $\rho_g = 0$ the exponent is exactly $1$, the relationship is linear, and we recover $\E[y_iy_j\mid\pi_{ij}] = V_A \pi_{ij}$.
Assortative mating shows up as a single number: a deficit in the exponent.

Second, the exponent is less than one, so the curve is concave.
The proportional inflation grows as pairs become more distant --- in the worked example the exponent is $0.823$, and while sibling covariance is inflated by $30\%$ relative to a randomly mating population with the same base heritability, first-cousin covariance is inflated by $66\%$.
This is the reason assortative mating is more visible in distant relatives than in close ones, and the reason it inflates estimates that lean on distant or nominally unrelated pairs.

Third, \eqref{eq:am-level2} is linear on a log--log scale,
$$ \log \E[\,y_iy_j \mid \pi_{ij}\,] \;=\; \log \VAeq \;+\; \big(1-\log_2(1+\rho_g)\big)\,\log \pi_{ij} , $$
so the slope of expected cross-product against realized relatedness, both on log scales, estimates $\rho_g$ through its shortfall below one.
Combined with the within-class result of \S\ref{sec:am-level23}, this gives two distinguishable signatures: \emph{within} a relationship class the response to sharing is linear with slope $\VAo$, while \emph{across} classes it is a power law with exponent below one.
Under random mating these are the same line.
The wedge between them is assortative mating.

Finally, note where \eqref{eq:am-level2} points as $\pi_{ij}\to 0$.
The covariance vanishes, but more slowly than linearly, and the pairs contributing in that limit are the nominally unrelated ones.
This is the regime in which the additive form \eqref{eq:am-level1} is the natural description, and where methods that fit a second, disequilibrium relatedness matrix alongside the GRM operate \citep{Zhang2025NatGenet}.

\subsection{Summary}
The same four information sets as Section 1, but no longer arranged in a chain.
Level 2 and Level 3 are parallel coarsenings of their common refinement, and the derivation runs $0 \to 1 \to 3 \to (2,3) \to 2$ rather than top to bottom.

\begin{center}
\small
\begin{tabular}{c l l p{4.0cm}}
\textbf{Level} & \textbf{Conditioned on} & $\E[y_i y_j \mid \cdot]$ & \textbf{New assumptions} \\
\hline
0 & causal genotypes, effects & $g_i g_j$ & \raggedright --- (unchanged) \tabularnewline
1 & genotypes $\bx_i,\bx_j$ & $\VAo A_{ij} + \text{disequilibrium}$ & \raggedright \assumptionid{Random-effects} weakened \tabularnewline
3 & pedigree & $\VAeq\left(\tfrac{1+\rho_g}{2}\right)^{d}$ (collateral) & \raggedright \assumptionid{Pheno-AM}, \assumptionid{Primary-AM}, \assumptionid{AM-equilibrium}, \assumptionid{Monogamy} \tabularnewline
$(2,3)$ & pedigree and $\pi_{ij}$ & $C_{ij} + \VAo(\pi_{ij}-\bar\pi_{d})$ & \raggedright \approximationid{Same-locus-AM}, \approximationid{GPD-pedigree} \tabularnewline
2 & $\pi_{ij}$ & $\VAeq\, \pi_{ij}^{\,1-\log_2(1+\rho_g)}$ & \raggedright \approximationid{Effective-degree} \tabularnewline
\hline
\end{tabular}
\end{center}

Reduced to one sentence: assortative mating adds a second channel of resemblance that travels through the pedigree rather than through shared segments, so relatedness and genealogy stop being interchangeable.
Everything else in this section is bookkeeping for that one fact.

\todo{Still to do in this section: the explicit disequilibrium relatedness matrix at Level 1 and its estimation in samples of unrelated individuals; the exact size of the residual term in \approximationid{GPD-pedigree}; and the non-equilibrium case, where a mismatch between observed and predicted correlations dates the onset of assortment \citep{Sunde2024NatComm}.}

\comment{For agents, the above is the end of the do-not-edit region opened earlier; everything from here down may be edited again, the bibliography included.}

\bibliographystyle{unsrtnat}
\bibliography{../../../bib/references}

\end{document}
